% Proofs for Section 4 (collectives).
\section{Proofs for \texorpdfstring{Section~\ref{sec:col}}{collectives}}\label{app:col}
% src: proofs transcribed from L2, L6, B1, B3, B5; gaps filled as noted in the comments after each proof

The proofs follow the order of \cref{sec:col}. Besides the lemmas below we use standard facts about symmetric and nonnegative matrices: the Courant--Fischer characterization of eigenvalues, the monotonicity of eigenvalues under addition of a positive semidefinite matrix, Cauchy interlacing, Gershgorin's theorem, the Perron--Frobenius theorem \citep{horn2012matrix}, and the Schur test, which bounds the operator norm of a block matrix with block norms $b_{ij}$ by $(\max_i\sum_jb_{ij}\cdot\max_j\sum_ib_{ij})^{1/2}$.

\begin{lemma}[Gaussian conditional dependence]\label{lem:col-gaussian}
Let $(U,Y,Z)$ be jointly Gaussian with covariance $\Sigma$ and precision $J=\Sigma^{-1}$, where $U$ is finite dimensional and $Y$ and $Z$ are finite or countable families, $\Sigma$ being a bounded operator on $\ell^2$ with bounded inverse.
\begin{enumerate}
\item[(a)] $U\perp Z\mid Y$ if and only if $J_{UZ}=0$.
\item[(b)] With $M=J_{UU}^{-1/2}J_{UZ}J_{ZZ}^{-1/2}$, $\Vert M\Vert<1$, $I(U;Z\mid Y)=-\frac12\ln\det(I-MM^\top)$, $\frac12\Vert M\Vert_F^2\le I(U;Z\mid Y)\le\Vert M\Vert_F^2/(2(1-\Vert M\Vert^2))$ and $\Vert M\Vert_F^2\le\lambda_{\max}(\Sigma)^2\Vert J_{UZ}\Vert_F^2$.
\end{enumerate}
\end{lemma}
% src: L2-L1, L6-L1; the countable case is added here for the infinite graphs of sec:col-markov
\begin{proof}
Conditional covariances of a Gaussian family are covariances of orthogonal projections in $L^2$, and for a bounded covariance with bounded inverse they are given by the block formulas of finite dimensions, with operator blocks of $J$ \citep{lauritzen1996graphical}. Given $(Y,Z)$, $U$ has covariance $J_{UU}^{-1}$; given $Y$ alone, $(U,Z)$ has precision $\begin{pmatrix}J_{UU}&J_{UZ}\\J_{ZU}&J_{ZZ}\end{pmatrix}$, so $U$ has covariance $(J_{UU}-J_{UZ}J_{ZZ}^{-1}J_{ZU})^{-1}$. These covariances do not depend on the conditioning values, and they agree exactly when the conditional mean of $U$ given $(Y,Z)$ does not depend on $Z$, that is when $U\perp Z\mid Y$; they agree exactly when $J_{UZ}J_{ZZ}^{-1}J_{ZU}=0$, that is when $J_{UZ}=0$, since $J_{ZZ}\succ0$. This is (a). For (b), with all determinants those of $\dim U\times\dim U$ matrices,
\[
I(U;Z\mid Y)=h(U\mid Y)-h(U\mid Y,Z)=\frac12\ln\frac{\det J_{UU}}{\det(J_{UU}-J_{UZ}J_{ZZ}^{-1}J_{ZU})}=-\frac12\ln\det(I-MM^\top),
\]
because $J_{UU}-J_{UZ}J_{ZZ}^{-1}J_{ZU}=J_{UU}^{1/2}(I-MM^\top)J_{UU}^{1/2}$. This matrix is positive definite, so $\Vert M\Vert<1$. The two bounds follow from $x\le-\ln(1-x)\le x/(1-x)$ on $[0,1)$, applied to the eigenvalues of $MM^\top$, each at most $\Vert M\Vert^2$. Principal blocks satisfy $\lambda_{\min}(J_{UU}),\lambda_{\min}(J_{ZZ})\ge\lambda_{\min}(J)=\lambda_{\max}(\Sigma)^{-1}$, which gives the last inequality.
\end{proof}

\begin{proof}[Proof of \cref{thm:col-path-blanket}]
(a) The path density is $p_{X_0}(x_0)\prod_{t<T}\Ncal(x_{t+1};Ax_t,Q^{\Xi})$, so the path is Gaussian with a precision $J$ that is block tridiagonal in time: $J_{00}=\Sigma_0^{-1}+A^\top(Q^{\Xi})^{-1}A$, $J_{tt}=(Q^{\Xi})^{-1}+A^\top(Q^{\Xi})^{-1}A$ for $1\le t\le T-1$, $J_{TT}=(Q^{\Xi})^{-1}$ and $J_{t+1,t}=-(Q^{\Xi})^{-1}A=J_{t,t+1}^\top$. By \cref{lem:col-gaussian}(a) the path blanket holds exactly when every block of $J$ between an interior coordinate at some time and an exterior coordinate at some time vanishes. The time-diagonal blocks give $(\Sigma_0^{-1}+A^\top(Q^{\Xi})^{-1}A)_{\sI\sE}=0$ at $t=0$, $((Q^{\Xi})^{-1}+A^\top(Q^{\Xi})^{-1}A)_{\sI\sE}=0$ at $1\le t\le T-1$, a nonempty range because $T\ge2$, and $((Q^{\Xi})^{-1})_{\sI\sE}=0$ at $t=T$. The lag-one blocks pair the interior at $t+1$ with the exterior at $t$ and the exterior at $t+1$ with the interior at $t$, giving $((Q^{\Xi})^{-1}A)_{\sI\sE}=0$ and $((Q^{\Xi})^{-1}A)_{\sE\sI}=0$; blocks at larger lags vanish. The three time-diagonal conditions are equivalent to the first, fourth and fifth conditions of \eqref{eq:col-pb-general}, and conversely \eqref{eq:col-pb-general} makes every listed block vanish. If $Q^{\Xi}$ is block diagonal for $(\sI,\sB,\sE)$, then $((Q^{\Xi})^{-1}A)_{\sI\sE}=(Q^{\Xi}_{\sI\sI})^{-1}A_{\sI\sE}$ and $((Q^{\Xi})^{-1}A)_{\sE\sI}=(Q^{\Xi}_{\sE\sE})^{-1}A_{\sE\sI}$ vanish exactly when $A_{\sI\sE}=0$ and $A_{\sE\sI}=0$, and then
\[
(A^\top(Q^{\Xi})^{-1}A)_{\sI\sE}=A_{\sI\sI}^\top(Q^{\Xi}_{\sI\sI})^{-1}A_{\sI\sE}+A_{\sB\sI}^\top(Q^{\Xi}_{\sB\sB})^{-1}A_{\sB\sE}+A_{\sE\sI}^\top(Q^{\Xi}_{\sE\sE})^{-1}A_{\sE\sE}=A_{\sB\sI}^\top(Q^{\Xi}_{\sB\sB})^{-1}A_{\sB\sE},
\]
while $(\Sigma_0^{-1})_{\sI\sE}=0$ is the initial conditional independence by \cref{lem:col-gaussian}(a). Sufficiency needs no Gaussian start. Let $\mu_0$ be any initial law with $X_{\sI,0}\perp X_{\sE,0}\mid X_{\sB,0}$, and let $R$ be the law of a path whose initial slice has law $\mu_0$ and whose later slices are independent of it, with independent standard normal coordinates. Under $R$ the interior and exterior paths are conditionally independent given the blanket path. The path law has density $\prod_{t<T}\Ncal(x_{t+1};Ax_t,Q^{\Xi})/\prod_{1\le t\le T}\Ncal(x_t;0,I)$ with respect to $R$, and under the first four conditions of \eqref{eq:col-pb-general} the exponent of the numerator has no term coupling an interior coordinate with an exterior one, so this density is a function of $(x_{\sI,0:T},x_{\sB,0:T})$ times a function of $(x_{\sE,0:T},x_{\sB,0:T})$. A density of this product form preserves the conditional independence, so the path blanket holds.

(b) By \cref{def:col-partition} no internal variable reads an exterior one and no exterior variable reads an internal one, so $A_{\sI\sE}=0$ and $A_{\sE\sI}=0$; active variables read no exterior variable, so $A_{k\sE}=0$ for $k\in\sA$. With per-channel noise $Q^{\Xi}_{\sB\sB}$ is block diagonal over channels, and the collider term reduces to $\sum_{k\in\sS}A_{k\sI}^\top(Q^{\Xi}_{kk})^{-1}A_{k\sE}$. A term is nonzero only if channel $k$ reads both an interior and an exterior variable. If a channel $k$ reads an interior variable $u$ and an exterior variable $u'$, the $(u,u')$ block of the sum contains $A_{ku}^\top(Q^{\Xi}_{kk})^{-1}A_{ku'}$, a polynomial in entries of $A$ that occur in no other term and that is not identically zero, because $(Q^{\Xi}_{kk})^{-1}\succ0$. The sum is therefore a polynomial in the coefficients that is not identically zero, and its zero set has positive codimension. With separate channels every sensory channel reads one source, which lies on one side, and the sum vanishes.

(c) Under the first three conditions of \eqref{eq:col-pb}, the sum $\sum_{t<T}\ln\Ncal(x_{t+1};Ax_t,Q^{\Xi})$ has no term that couples an interior coordinate with an exterior one, so it equals $f(x_{\sI,0:T},x_{\sB,0:T})+g(x_{\sE,0:T},x_{\sB,0:T})$. Fix a blanket path $x_{\sB}$ and write $Y_\sI=X_{\sI,0:T}$, $Y_\sE=X_{\sE,0:T}$. Their conditional density is proportional to $p_{X_0}(y_\sI(0),x_\sB(0),y_\sE(0))e^{f}e^{g}$, a function of $(y_\sI(0),y_\sE(0))$ times a function of $y_\sI$ times a function of $y_\sE$. Hence, given $x_\sB$, $Y_\sI(1{:}T)\perp Y_\sE\mid Y_\sI(0)$ and $Y_\sE(1{:}T)\perp Y_\sI\mid Y_\sE(0)$, and the chain rule gives
\begin{align*}
I(Y_\sI;Y_\sE\mid x_\sB)&=I(Y_\sI(0);Y_\sE\mid x_\sB)=I(Y_\sI(0);Y_\sE(0)\mid x_\sB)+I(Y_\sI(0);Y_\sE(1{:}T)\mid Y_\sE(0),x_\sB)\\
&=I(Y_\sI(0);Y_\sE(0)\mid x_\sB).
\end{align*}
Averaging over $x_\sB$ gives the initial-slice formula. For the bound, $I(U;Z\mid Y)\le I(U;Y,Z)$ and the chain rule give
\[
\iota_T\le I(X_{\sI,0};X_{\sB,0:T},X_{\sE,0})=I(X_{\sI,0};X_{\sB,0},X_{\sE,0})+I(X_{\sI,0};X_{\sB,1:T}\mid X_{\sB,0},X_{\sE,0}).
\]
Write $X_{\sB,1:T}=\mathcal LX_0+\tilde\Xi_{\sB}$, where $\mathcal L$ stacks the blanket rows of $A^t$, $t=1,\dots,T$, and $\tilde\Xi_{\sB}=\mathcal J_\sB\Xi_{0:T-1}$ is the response of the blanket to the noise, independent of $X_0$. Given $(X_{\sB,0},X_{\sE,0})$, $X_{\sI,0}$ is Gaussian with covariance at most $\Vert\Sigma_0\Vert I$, and the part of $\mathcal LX_0$ not driven by $X_{\sI,0}$ is known, so the Gaussian channel formula and $\ln\det(I+X)\le\tr X$ give
\[
I(X_{\sI,0};X_{\sB,1:T}\mid X_{\sB,0},X_{\sE,0})\le\tfrac12\Vert\Sigma_0\Vert\Vert\Cov(\tilde\Xi_{\sB})^{-1}\Vert\Vert\mathcal L\Vert_F^2.
\]
The map $\mathcal J$ from $\Xi_{0:T-1}$ to $X_{1:T}$ at zero initial state is block lower triangular with identity diagonal blocks, and its inverse is block bidiagonal with blocks $I$ and $-A$, so its smallest singular value is at least $(1+\Vert A\Vert)^{-1}$. Since $\mathcal J_\sB\mathcal J_\sB^\top$ is a principal submatrix of $\mathcal J\mathcal J^\top$, Cauchy interlacing gives $\lambda_{\min}(\mathcal J_\sB\mathcal J_\sB^\top)\ge(1+\Vert A\Vert)^{-2}$, and $\Cov(\tilde\Xi_{\sB})=\mathcal J_\sB(I\otimes Q^{\Xi})\mathcal J_\sB^\top\succeq\lambda_{\min}(Q^{\Xi})\mathcal J_\sB\mathcal J_\sB^\top$. Finally $\Vert\mathcal L\Vert_F^2\le\sum_{t\ge1}\Vert A^t\Vert_F^2$, finite when $\rho(A)<1$, which gives \eqref{eq:col-initial-slice}.
\end{proof}
% src: L2-T1 proof, L2-P1 proof, L6-P1(a),(c) proof; the positive-codimension step of (b) is written out here

\begin{proof}[Proof of \cref{ex:col-collider}]
Here $A_{\sB\sI}^\top(Q^{\Xi}_{\sB\sB})^{-1}A_{\sB\sE}=\alpha_1\alpha_2\neq0$. For $t\ge1$ the triple formed by $X_{i,t}$, $X_{e,t}$ and $X_{b,t+1}$ is a function of the noises $\Xi^i_{t-1}$, $\Xi^e_{t-1}$ and $\Xi^b_t$, so for every initial law these triples, for $1\le t\le T-1$, are independent of one another and of the variables of rounds $0$, $1$ and $T$ that they do not contain. The window leak is therefore the sum of a term that involves only round $0$, the initial state and $X_{b,1}$, and of $T-1$ equal one-round terms $I(X_{i,t};X_{e,t}\mid X_{b,t+1})$; on every window with $T\ge2$ it is positive, and the leak rate is the one-round term whenever the round-$0$ term is finite. In one round, $X_{i,t}$ and $X_{e,t}$ are independent standard normal and $X_{b,t+1}=\alpha_1X_{i,t}+\alpha_2X_{e,t}+\Xi^b_t$, so their conditional covariance given $X_{b,t+1}$ is $I-vv^\top/(1+\alpha_1^2+\alpha_2^2)$ with $v=(\alpha_1,\alpha_2)$. The squared partial correlation is $\alpha_1^2\alpha_2^2/((1+\alpha_1^2)(1+\alpha_2^2))$, and since $(1+\alpha_1^2)(1+\alpha_2^2)-\alpha_1^2\alpha_2^2=1+\alpha_1^2+\alpha_2^2$, the conditional mutual information $-\frac12\ln(1-\alpha_1^2\alpha_2^2/((1+\alpha_1^2)(1+\alpha_2^2)))$ equals \eqref{eq:col-collider}.
\end{proof}
% src: L2-X1, C_critic V1; the decomposition for an arbitrary initial law is written out here

\begin{proof}[Proof of \cref{prop:col-two-hop}]
If $d(\sI,\sE)\ge2$, one-hop coupling gives $A_{\sI\sE}=0$ and $A_{\sE\sI}=0$. A nonzero term $A_{k\sI}^\top(Q^{\Xi}_{kk})^{-1}A_{k\sE}$ requires an agent $k$ within one hop of some $i\in\sI$ and of some $e\in\sE$, hence $d(i,e)\le2$. When $d(\sI,\sE)\ge3$ every term vanishes, and \cref{thm:col-path-blanket}(a) gives exactness; the stationary start is Gaussian with covariance $\sum_nA^nQ^{\Xi}(A^\top)^n\succeq Q^{\Xi}\succ0$, so \cref{thm:col-path-blanket}(c) gives a zero leak rate. When $d(\sI,\sE)=2$ the collider term is the stated sum, and a shell agent that reads $i\in\sI$ and $e\in\sE$ contributes to the $(i,e)$ block a term in entries of $A$ that no other agent's term contains, so the argument of \cref{thm:col-path-blanket}(b) applies. For a schedule $A_t$ the path precision has the blocks $\Sigma_0^{-1}+A_0^\top(Q^{\Xi})^{-1}A_0$, $(Q^{\Xi})^{-1}+A_t^\top(Q^{\Xi})^{-1}A_t$ and $-(Q^{\Xi})^{-1}A_t$ with the same sparsity, and the proof applies round by round. The shell $\{k:1\le d(k,\sI)\le2\}$ makes $d(\sI,\sE)=3$, and every subset $\sI$ of the agents gives one such partition.
\end{proof}
% src: L2-P6 proof, L6-P1 corollary

\begin{proof}[Proof of \cref{prop:col-records-partition}]
(a) Per-channel noise is independent across the partition, so by \cref{thm:col-path-blanket}(a) we check $A_{\sI\sE}=0$, $A_{\sE\sI}=0$ and the collider term. An internal record reads only its member writers and is read only by member channels, so it reads and drives no exterior variable. A boundary record reads only its writers, which are collective, and lies in $\sA$, which may drive exterior variables. A member channel that writes only internal or boundary records and reads nothing outside drives and reads only collective variables, so it lies in $\sI$ and has no exterior neighbor. An exterior record is written by member channels, which lie in $\sA$ or $\sS$, and read by member channels that read nothing else, which lie in $\sS$ and read one variable. Hence $A_{\sI\sE}=A_{\sE\sI}=0$, record-reading channels contribute no collider term, and the collider condition is that of the remaining sensory channels, which is the condition of \cref{thm:col-path-blanket}(b). The persistence of a record enters only the record's own row of $A$, on which no condition bears, and records are state variables of the one-step Markov process \eqref{eq:col-joint}, so lags through them need no further condition.

(b) A record read by an outsider and placed in $\sI$ makes an entry of $A_{\sE\sI}$ nonzero. A record with writers on both sides placed in $\sB$ makes its member writers, which then read nothing outside and drive only collective variables, internal, while its other writers are exterior; the record reads an interior and an exterior variable in one round, and its collider term $A_{r\sI}^\top(Q^{\Xi}_{rr})^{-1}A_{r\sE}$ is nonzero for generic coefficients. An exterior record whose writers are all non-members, placed in $\sS$, reads only exterior variables and itself, contributes no collider term, and leaves the conditions of \cref{thm:col-path-blanket}(a) unchanged.
\end{proof}
% src: B5-P6 proof

\begin{proof}[Proof of \cref{prop:col-memory}]
(a) The path blanket on $[0,t]$ marginalizes to $X_{\sI,t}\perp X_{\sE,t}\mid X_{\sB,0:t}$. By the chain rule,
\[
I(X_{\sI,t};X_{\sE,t}\mid X_{\sB,t})\le I(X_{\sI,t};X_{\sE,t},X_{\sB,0:t-1}\mid X_{\sB,t})=I(X_{\sI,t};X_{\sB,0:t-1}\mid X_{\sB,t})+I(X_{\sI,t};X_{\sE,t}\mid X_{\sB,0:t}),
\]
and the last term vanishes; exchanging interior and exterior gives the other bound. (b) By (a) and Gaussianity, the conditional density of $(X_{\sI,t},X_{\sE,t})$ given the history $H=X_{\sB,0:t}$ is $p(x_\sI\mid H)\Ncal(x_\sE;\hat X_{\sE,t}(H),C_\sE)$ with a covariance $C_\sE$ that does not depend on $H$. Conditioning on $(X_{\sB,t},\hat X_{\sE,t})$, a function of $H$, instead of on $H$ averages the first factor over $H$ and leaves the second unchanged, because it depends on $H$ only through $\hat X_{\sE,t}$; the result is a function of $x_\sI$ times a function of $x_\sE$. (c) From $AQ^{\Xi}=Q^{\Xi}A^\top$, induction gives $A^kQ^{\Xi}(A^\top)^k=A^{2k}Q^{\Xi}$, so $\Sigma=(I-A^2)^{-1}Q^{\Xi}$ and $A\Sigma=(I-A^2)^{-1}AQ^{\Xi}=\Sigma A^\top$, which is time reversibility of the stationary chain. Multiplying $AQ^{\Xi}=Q^{\Xi}A^\top$ by $(Q^{\Xi})^{-1}$ on both sides gives $(Q^{\Xi})^{-1}A=A^\top(Q^{\Xi})^{-1}$, hence $(Q^{\Xi})^{-1}A^2=A^\top(Q^{\Xi})^{-1}A$ and $\Pi=(Q^{\Xi})^{-1}(I-A^2)=(Q^{\Xi})^{-1}-A^\top(Q^{\Xi})^{-1}A$. Under the block conditions the expansion in the proof of \cref{thm:col-path-blanket}(a) gives $\Pi_{\sI\sE}=-A_{\sB\sI}^\top(Q^{\Xi}_{\sB\sB})^{-1}A_{\sB\sE}$, and \cref{lem:col-gaussian}(a) finishes the proof.
\end{proof}
% src: L2-P5 proof, L2-P4 proof

\begin{proof}[Proof of \cref{thm:col-thick-shell}]
(a) Fix a vertex $j_0$ and $R>0$, let $f(i)=\min\{d(i,j_0),R\}$, which is $1$-Lipschitz and bounded, and let $\mathcal T=\diag(e^{m_\Pi f(i)}I)$ act blockwise; $\mathcal T$ and $\mathcal T^{-1}$ are bounded. The matrix $\mathcal T\Sigma\mathcal T^{-1}-\Sigma$ has blocks $\Sigma_{ij}(e^{m_\Pi(f(i)-f(j))}-1)$, of norm at most $B_\Sigma e^{-m_\Sigma d(i,j)}(e^{m_\Pi d(i,j)}-1)$ because $|e^x-1|\le e^{|x|}-1$ and $|f(i)-f(j)|\le d(i,j)$. Its block row and column sums are at most $\mathcal C(m_\Pi)$, so the Schur test bounds its norm by $\mathcal C(m_\Pi)$. Then $\mathcal T\Sigma\mathcal T^{-1}=\Sigma(I+\Sigma^{-1}(\mathcal T\Sigma\mathcal T^{-1}-\Sigma))$ with $\Vert\Sigma^{-1}(\mathcal T\Sigma\mathcal T^{-1}-\Sigma)\Vert\le\mathcal C(m_\Pi)/\lambda_{\min}(\Sigma)<1$, so $\Vert(\mathcal T\Sigma\mathcal T^{-1})^{-1}\Vert\le(\lambda_{\min}(\Sigma)-\mathcal C(m_\Pi))^{-1}=B_\Pi$. The $(i,j_0)$ block of $(\mathcal T\Sigma\mathcal T^{-1})^{-1}=\mathcal T\Pi\mathcal T^{-1}$ is $e^{m_\Pi f(i)}\Pi_{ij_0}$, of norm at most $B_\Pi$, and letting $R\to\infty$ gives the claim. (b) Since $\sI$ is finite, $e_\sI$ is finite dimensional, and \cref{lem:col-gaussian}(b) applies with $U=e_\sI$, $Y=e_\sB$, $Z=e_\sE$: $\Vert M\Vert_F^2\le\lambda_{\max}(\Sigma)^2\Vert\Pi_{\sI\sE}\Vert_F^2\le\lambda_{\max}(\Sigma)^2Z_{\sI\sE}$, since a block of size at most $K$ has $\Vert\Pi_{ik}\Vert_F^2\le K\Vert\Pi_{ik}\Vert^2$. With $\sigma_l$ the singular values of $M$, $\sigma_l^2\le\Vert M\Vert_F^2<1$, so $-\ln(1-\sigma_l^2)\le\sigma_l^2/(1-\Vert M\Vert_F^2)$, and summing gives \eqref{eq:col-thick-shell}. (c) For $i\in\sI$, $\sum_{k\in\sE}e^{-2m_\Pi d(i,k)}\le\Psi(d(i,\sE))$ with $\Psi(t)=\sum_{r\ge t}\mathcal S(r)e^{-2m_\Pi r}$, and
\begin{align*}
\sum_{i\in\sI}\Psi(d(i,\sE))&\le n_\partial\sum_{t\ge h}\Psi(t)=n_\partial\sum_{r\ge h}(r-h+1)\mathcal S(r)e^{-2m_\Pi r}\\
&\le n_\partial B_{\mathcal G}(1+h)^{D-1}e^{-2m_\Pi h}\sum_{j\ge0}(1+j)^De^{-2m_\Pi j},
\end{align*}
writing $r=h+j$ and using $h+j\le(1+h)(1+j)$. Take $B_*=B_{\mathcal G}\sum_{j\ge0}(1+j)^De^{-2m_\Pi j}$.
\end{proof}
% src: L2-L2 proof (truncation of the weight added here for infinite graphs), L2-T2 proof

\begin{proof}[Proof of \cref{thm:col-markov-bound}]
(a) Fix a vertex $i$ and a width $h$ and apply \cref{lem:col-gaussian}(b) with $U=e_i$, $Y=e_{\sB_h(i)}$ and $Z=e_{\sE_h(i)}$. Then
\[
\Vert M\Vert_F^2\le\lambda_{\max}(\Sigma)^2K\sum_{j\in\sE_h(i)}\Vert\Pi_{ij}\Vert^2\le\lambda_{\max}(\Sigma)^2KB_\Pi^2B_{\mathcal G}\sum_{r>h}r^{D-1}e^{-2m_\Pi r},
\]
uniformly in $i$. For every $\xi'>1/(2m_\Pi)$ the right side is at most a constant times $e^{-h/\xi'}$, and once it is at most $\frac14$, $\iota(i,h)\le\frac23\Vert M\Vert_F^2$. For the finitely many smaller $h$ we use a bound uniform in $i$: $J_{UU}-J_{UZ}J_{ZZ}^{-1}J_{ZU}$ is the inverse of a conditional covariance, hence at least $\lambda_{\max}(\Sigma)^{-1}I$, while $J_{UU}\preceq\lambda_{\min}(\Sigma)^{-1}I$, so $I-MM^\top=J_{UU}^{-1/2}(J_{UU}-J_{UZ}J_{ZZ}^{-1}J_{ZU})J_{UU}^{-1/2}\succeq(\lambda_{\min}(\Sigma)/\lambda_{\max}(\Sigma))I$ and $\iota(i,h)\le\frac K2\ln(\lambda_{\max}(\Sigma)/\lambda_{\min}(\Sigma))$. Hence $\sup_i\iota(i,h)\le B'e^{-h/\xi'}$ for a constant $B'$ and all $h$, and $\xi_M\le1/(2m_\Pi)$.

(b) Let $0<m_\Pi\le m_\Sigma/(D+1)$ and $m'=m_\Sigma-m_\Pi\ge Dm_\Sigma/(D+1)$. Since $e^x-1\le xe^x$ for $x\ge0$, $\mathcal C(m_\Pi)\le B_\Sigma B_{\mathcal G}m_\Pi\sum_{r\ge1}r^De^{-m'r}$. For $x\in[r,r+1]$, $r^De^{-m'r}\le e^{m'}x^De^{-m'x}$, so $\sum_{r\ge1}r^De^{-m'r}\le e^{m'}\int_0^\infty x^De^{-m'x}\mathrm{d}x=e^{m'}D!m'^{-(D+1)}$. With $e^{m'}\le e^{m_\Sigma}$ and $m'^{-(D+1)}\le(1+1/D)^{D+1}m_\Sigma^{-(D+1)}\le4m_\Sigma^{-(D+1)}$, because $(1+1/D)^{D+1}\le4$ for $D\ge1$,
\[
\mathcal C(m_\Pi)\le4B_\Sigma B_{\mathcal G}D!e^{m_\Sigma}m_\Pi m_\Sigma^{-(D+1)}.
\]
The choice $m_\Pi=m_*$ of \eqref{eq:col-mstar} satisfies both constraints and makes $\mathcal C(m_*)\le\lambda_{\min}(\Sigma)/2$, so \cref{thm:col-thick-shell}(a) gives $B_\Pi\le2/\lambda_{\min}(\Sigma)$, and (a) gives \eqref{eq:col-markov-bound}.

(c) The hypotheses are those of \cref{thm:sca-clustering}(b) and of its last clause: a time-invariant $W$ satisfies \textup{(Lin)}, the law of $\zeta(t)$ is time invariant by hypothesis, so that its covariance $Q^{\zeta}$ does not depend on $t$, and \textup{(B)} gives block-diagonal local noise with blocks between $q_{\min}I$ and $qI$. We repeat the argument for the stationary covariance so that the constants are explicit. The covariance is $\Sigma=\sum_{n\ge0}W^nQ^{\zeta}(W^\top)^n$ by \eqref{eq:col-sigma-c}. By \textup{(Loc)}, $(W^n)_{ik}=0$ when $d(i,k)>n$, and by \textup{(C$_r$)} and induction $\sum_k\Vert(W^n)_{ik}\Vert\le\lambda^n$. Block $(i,j)$ of the $n$th term is $\sum_k(W^n)_{ik}Q^{\zeta}_{kk}(W^n)_{jk}^\top$, which vanishes unless $2n\ge d(i,j)$ and otherwise has norm at most $q\max_k\Vert(W^n)_{jk}\Vert\sum_k\Vert(W^n)_{ik}\Vert\le q\lambda^{2n}$. Summing over $n\ge\lceil d(i,j)/2\rceil$ gives $\Vert\Sigma_{ij}\Vert\le q\lambda^{2\lceil d(i,j)/2\rceil}/(1-\lambda^2)\le q\lambda^{d(i,j)}/(1-\lambda^2)$, which is \eqref{eq:sca-clustering}, so $B_\Sigma=q/(1-\lambda^2)$, $m_\Sigma=\ln(1/\lambda)$ and $e^{m_\Sigma}=1/\lambda$. Every term of the series is positive semidefinite and the term $n=0$ is $Q^{\zeta}$, so $\Sigma\succeq Q^{\zeta}\succeq q_{\min}I$, the last clause of \cref{thm:sca-clustering}, and the Schur test gives $\lambda_{\max}(\Sigma)\le B_\Sigma\sum_{r\ge0}\mathcal S(r)\lambda^r<\infty$, with $\mathcal S(0)=1$. Replacing $\lambda_{\min}(\Sigma)$ by its lower bound $q_{\min}$ in \eqref{eq:col-mstar} lowers $m_*$ and keeps $\mathcal C(m_*)\le q_{\min}/2\le\lambda_{\min}(\Sigma)/2$, which gives the stated bound.
\end{proof}
% src: B3-T2 proof, L6-P3 proof (uniform bound for small widths added here); (c) re-derives the constants of thm:sca-clustering(b) so that the dependence is explicit

\begin{remark}[A dynamical route to \cref{conj:col-markov-length}]\label{rem:col-backward}
From $\Sigma=W\Sigma W^\top+Q^{\zeta}$ and the Woodbury identity, $\Pi=(Q^{\zeta})^{-1}-(Q^{\zeta})^{-1}W\Sigma_{\leftarrow}W^\top(Q^{\zeta})^{-1}$ with $\Sigma_{\leftarrow}=(\Pi+W^\top(Q^{\zeta})^{-1}W)^{-1}=\Cov(e(t-1)\mid e(t))$: the long-range part of the precision is the conditional covariance of the previous state given the present one, pushed through one hop on each side. Since $0\preceq\Pi\preceq(Q^{\zeta})^{-1}$, the inverse of $\Sigma_{\leftarrow}$ lies between $\sigma_{\min}(W)^2/q$ and $(1+\Vert W\Vert^2)/q_{\min}$, a condition number free of $\lambda_{\max}(\Sigma)$, which is what diverges in the massless limit; for normal $W$ with $Q^{\zeta}=qI$, $\Sigma_{\leftarrow}=qI$ and the range-two precision of \cref{thm:col-diffusive-precision}(b) follows. A proof of \cref{conj:col-markov-length} along this route needs decay of $\Sigma_{\leftarrow}$ from decay of $\Pi$ without the loss of a Combes--Thomas step.
\end{remark}
% src: B3-R2

\begin{proof}[Proof of \cref{prop:col-markov-vs-correlation}]
(a) In the stationary state $e^{\mathrm{src}}_k(t)=\zeta^{\mathrm{src}}_k(t-1)$ and $e^{\mathrm{lis}}_k(t)=a[\zeta^{\mathrm{src}}_{k-1}(t-2)+\zeta^{\mathrm{src}}_k(t-2)]+\zeta^{\mathrm{lis}}_k(t-1)$. At equal times sources are independent of one another and of every listener; listeners $k$ and $k+1$, two hops apart, share $\zeta^{\mathrm{src}}_k(t-2)$, so $\Cov(e^{\mathrm{lis}}_k,e^{\mathrm{lis}}_{k+1})=a^2q_{\mathrm{src}}$ and $\Var(e^{\mathrm{lis}}_k)=2a^2q_{\mathrm{src}}+q_{\mathrm{lis}}$; all other covariances vanish. The precision is block diagonal over sources and listeners, with $q_{\mathrm{src}}^{-1}I$ on sources. The listener covariance is the Laurent operator with diagonal $\sigma_0^2=2a^2q_{\mathrm{src}}+q_{\mathrm{lis}}$ and off-diagonal $\sigma_1^2=a^2q_{\mathrm{src}}$, whose symbol $p\mapsto\sigma_0^2+2\sigma_1^2\cos p$ has range $[q_{\mathrm{lis}},4a^2q_{\mathrm{src}}+q_{\mathrm{lis}}]$. Its inverse $\Pi^{\mathrm{lis}}$ has entries
\[
\Pi^{\mathrm{lis}}_{k,k+n}=\frac1{2\pi}\int_{-\pi}^{\pi}\frac{e^{inp}}{\sigma_0^2+2\sigma_1^2\cos p}\mathrm{d}p=\frac{(-1)^ne^{-m_\psi|n|}}{\sqrt{\sigma_0^4-4\sigma_1^4}},\qquad2\cosh m_\psi=\frac{\sigma_0^2}{\sigma_1^2}=2+\psi,
\]
by the residue at the root $-e^{-m_\psi}\in(-1,0)$ of $\sigma_1^2x^2+\sigma_0^2x+\sigma_1^2$, so $m_\psi=\arccosh(1+\psi/2)$ and $\sqrt{\sigma_0^4-4\sigma_1^4}=a^2q_{\mathrm{src}}\sqrt{\psi(4+\psi)}$. Sources drop out of every conditional mutual information, and a source vertex has $\iota=0$. For listener $0$ and width $h$ the exterior listeners are those with $|k|\ge n_h=\lfloor h/2\rfloor+1$, and \cref{lem:col-gaussian}(b) applies to the listener precision. Every compression of $\Pi^{\mathrm{lis}}$ to a set of listeners has spectrum in $[(4a^2q_{\mathrm{src}}+q_{\mathrm{lis}})^{-1},q_{\mathrm{lis}}^{-1}]$, so $\Vert M\Vert_F^2=(\Pi^{\mathrm{lis}}_{00})^{-1}\Pi^{\mathrm{lis}}_{0\sE}(\Pi^{\mathrm{lis}}_{\sE\sE})^{-1}\Pi^{\mathrm{lis}}_{\sE0}$ lies between constants times $\sum_{|n|\ge n_h}e^{-2m_\psi|n|}$, which is of order $e^{-2m_\psi n_h}$. Since $\Vert M\Vert\to0$, $\iota(0,h)$ is bounded above and below by constants times $e^{-2m_\psi n_h}$, and $h+1\le2n_h\le h+2$ gives $\xi_M=1/m_\psi=1/\arccosh(1+\psi/2)$. The expansion $\arccosh(1+x)=\sqrt{2x}(1+O(x))$ gives the asymptotics. The spectrum of $\Sigma$ is $\{q_{\mathrm{src}}\}\cup[q_{\mathrm{lis}},4a^2q_{\mathrm{src}}+q_{\mathrm{lis}}]$, so $\operatorname{cond}(\Sigma)=q_{\mathrm{src}}/q_{\mathrm{lis}}=1/(a^2\psi)$ when $4a^2q_{\mathrm{src}}+q_{\mathrm{lis}}\le q_{\mathrm{src}}$, that is when $\psi\le(1-4a^2)/a^2$. Otherwise the largest eigenvalue is $4a^2q_{\mathrm{src}}+q_{\mathrm{lis}}$, and $\operatorname{cond}(\Sigma)=1+4/\psi$ as long as $q_{\mathrm{lis}}\le q_{\mathrm{src}}$.

(b) The covariance of $h+2$ consecutive sites is the tridiagonal Toeplitz matrix $T_{h+2}$ with diagonal $1+\alpha^2$ and off-diagonal $\alpha$, so covariances vanish beyond one site. For a tridiagonal matrix the cofactor formula gives $|(T_n^{-1})_{1n}|=\alpha^{n-1}/\det T_n$ and $(T_n^{-1})_{11}=(T_n^{-1})_{nn}=\det T_{n-1}/\det T_n$, and the recursion $\det T_k=(1+\alpha^2)\det T_{k-1}-\alpha^2\det T_{k-2}$ gives $\det T_k=\sum_{j=0}^k\alpha^{2j}$. The partial correlation of the end sites is therefore $\alpha^{n-1}/\det T_{n-1}$, which with $n=h+2$ is the stated value. The conditional mutual information of the end sites is $-\frac12\ln(1-\mathrm{pc}_h^2)\ge\frac12\mathrm{pc}_h^2$, with $\mathrm{pc}_h$ this partial correlation, and the half-line quantity is at least as large by the chain rule. For $\alpha<1$, $\mathrm{pc}_h^2\ge(1-\alpha^2)^2\alpha^{2h+2}$, so no bound $B'e^{-h/\xi'}$ holds with $\xi'<1/(2\ln(1/\alpha))$; at $\alpha=1$ the quantity is at least $\frac12(h+2)^{-2}$, which decays slower than every exponential.

(c) The update matrix is symmetric, hence normal, and $Q^{\zeta}=qI$, so \cref{thm:col-diffusive-precision}(c) gives a precision of range two; $\Pi^c_{i,\sE_h(i)}=0$ for $h\ge2$, and \cref{lem:col-gaussian}(a) gives $\iota(i,h)=0$, so $\xi_M=0$. The divergence of the Ornstein--Zernike length as $\varepsilon\to0$ is \cref{thm:sca-oz}(c).
\end{proof}
% src: B3-P5 proof, L6-P4 proof, C_critic V2, C2-7

\begin{proof}[Proof of \cref{prop:col-selects-nothing}]
Part (a) is \cref{thm:col-path-blanket}(b). In (b), path blankets are \cref{prop:col-two-hop}(a). The stationary precision is $(Q^{\Xi})^{-1}-A^\top(Q^{\Xi})^{-1}A$ for reversible dynamics (\cref{prop:col-memory}(c)) and $(I-AA^\top)/q$ for normal dynamics with $Q^{\Xi}=qI$ (as in \cref{thm:col-diffusive-precision}(b), with $A$ for $W$ and $Q^{\Xi}$ for $Q^\zeta$); by the argument of \cref{thm:col-diffusive-precision}(c) a block $(i,k)$ of either vanishes unless some agent is within one hop of both, so the interior-exterior block vanishes when $d(\sI,\sE)\ge3$, and \cref{lem:col-gaussian}(a) gives the stationary blanket. Part (c) is \cref{thm:col-thick-shell}(c), and (d) follows from the translation invariance of the dynamics and of the stationary law. Under (a) and (b) every candidate of the stated classes thus has an exact blanket and zero path leak; under (c) the leak bound, and with it the tolerance that a thick enough shell certifies, depends on the region only through $n_\partial$ and $h$; and by (d) neither blankets nor leaks separate translates. In the reversible and normal isotropic cases the instantaneous leak of every region with a two-hop shell vanishes as well, which makes leak minimization degenerate.
\end{proof}
% src: L2-P7 proof, L2-P8 proof

\begin{proof}[Proof of \cref{prop:col-world-conditional}]
In the level-0 network of \cref{def:fw-level0} every channel noise is specific to its receiver, so \cref{def:col-world-conditional} conditions on the world path alone, and the world-conditional law is the conditional law given the world path. By \cref{lem:fw-common-mode}(a) and (c) this law is Gaussian, with mean $e^{\mathrm{com}}$ and the law of $e^{\mathrm{loc}}$ around it, and by \cref{lem:fw-common-mode}(b) its covariance at one time is $\Sigma^{\mathrm{loc}}$, which is \eqref{eq:col-sigma-c} and does not depend on the path. A sender-side noise shared by several receivers adds to row $i$ of \eqref{eq:fw-level0-error} the term $F\Kg_{ij}$ times that noise for every attended sender $j$, a term of the same kind as the world innovation: it is independent of the world and of the local noises. The argument of \cref{lem:fw-common-mode}(b) and (c), run with the $\sigma$-algebra generated by the world path and the shared noises in place of that of the world path, puts the response of the network to these terms, a series in powers of $W$ like \eqref{eq:fw-common-mode}, into the common part, which stays measurable for the enlarged conditioning, and leaves the connected part and its law unchanged.
\end{proof}
% src: B1-P1 proof, through lem:fw-common-mode

\begin{proof}[Proof of \cref{thm:col-diffusive-precision}]
(a) Induction on $WQ^{\zeta}=Q^{\zeta}W^\top$ gives $W^nQ^{\zeta}=Q^{\zeta}(W^\top)^n$ for every $n$, hence
\[
W^nQ^{\zeta}(W^\top)^n=W^{2n}Q^{\zeta},\qquad\Sigma^c=\sum_{n\ge0}W^{2n}Q^{\zeta}=(I-W^2)^{-1}Q^{\zeta}.
\]
As in the proof of \cref{prop:col-memory}(c), $(Q^{\zeta})^{-1}W=W^\top(Q^{\zeta})^{-1}$ and $\Pi^c=(Q^{\zeta})^{-1}-W^\top(Q^{\zeta})^{-1}W$, for block matrices as for scalar ones. For scalar errors the off-diagonal entries are $-\sum_kW_{ki}W_{kj}/q_k$ and the sum of row $i$ is $1/q_i-\sum_kW_{ki}\lambda_k/q_k$; the identity $(Q^{\zeta})^{-1}W=W^\top(Q^{\zeta})^{-1}$ reads $W_{ki}/q_k=W_{ik}/q_i$, which turns the row sum into $(1-\sum_kW_{ik}\lambda_k)/q_i$. If $W\ge0$ with row sums at most $\lambda$, the off-diagonal entries are nonpositive and $\sum_kW_{ik}\lambda_k\le\lambda_i\max_k\lambda_k\le\lambda^2$. (b) Normality gives $W^n(W^\top)^n=(WW^\top)^n$, and $\Vert WW^\top\Vert=\Vert W\Vert^2=\rho(W)^2<1$ for normal $W$, so $\Sigma^c=q\sum_n(WW^\top)^n=q(I-WW^\top)^{-1}$. For scalar errors the off-diagonal entries are $-q^{-1}\sum_kW_{ik}W_{jk}$ and the sum of row $i$ is $(1-(WW^\top\one)_i)/q=(1-\sum_kW_{ik}\lambda^{\mathrm{col}}_k)/q$; with $W\ge0$ diffusivity is nonnegativity of these sums, and $\sum_kW_{ik}\lambda^{\mathrm{col}}_k\le\lambda_i\max_k\lambda^{\mathrm{col}}_k$. (c) With $Q^{\zeta}$ block diagonal, block $(i,j)$ of $W^\top(Q^{\zeta})^{-1}W$ is $\sum_kW_{ki}^\top(Q^{\zeta}_{kk})^{-1}W_{kj}$ and block $(i,j)$ of $WW^\top$ is $\sum_kW_{ik}W_{jk}^\top$. A term of the first is nonzero only if $k$ attends both $i$ and $j$, and a term of the second only if $i$ and $j$ both attend $k$, self-weights included; either way $d(i,j)\le2$. When $d(\sI,\sE)\ge3$, $\Pi^c_{\sI\sE}=0$, and \cref{lem:col-gaussian}(a) gives the stationary blanket. (d) The sum of row $i$ is $F(1-k^s_i-\sum_jk_{ij})+F\sum_jk_{ij}=F(1-k^s_i)$, and reversibility $W_{ij}q_j=q_iW_{ji}$ is $k_{ij}q_j=q_ik_{ji}$. With a common sensing gain every row sum equals $\lambda$, so $\sum_kW_{ik}\lambda_k=\lambda^2$ in case (a); in case (b), $W\one=\lambda\one$ gives $W^\top\one=\lambda\one$, because a real normal matrix and its transpose have the same eigenvectors with conjugate eigenvalues, so every column sum is $\lambda$ as well. (e) $W$ is symmetric with $Q^{\zeta}=qI$, so (b) gives $\Pi^c=(I-W^2)/q$. In index form $(W^2)_{ij}=\sum_kW_{ik}W_{kj}$: for neighbors $i,j$ only $k\in\{i,j\}$ contributes, giving $2ca$; for sites two steps apart on one axis only the midpoint contributes, giving $a^2$; for sites displaced by one step in each of two coordinates the two common neighbors contribute, giving $2a^2$; and every row of $W^2$ sums to $\lambda^2$, which gives the leak.
\end{proof}
% src: B1-T1 proof, L2-R1, L6-P2; (c) for vector errors and the column sums of (d) written out here

\begin{proof}[Proof of \cref{ex:col-non-diffusive}(a)]
Agents $1$ and $2$ are driven by independent noises, so $\Sigma^c_{12}=0$. Stationarity gives $e_3(t)=\sum_{n\ge0}W_{33}^n(W_{31}e_1(t-1-n)+W_{32}e_2(t-1-n)+\zeta_3(t-1-n))$, and since $\Cov(e_1(t),e_1(t-1-n))=W_{11}^{n+1}\Var(e_1)$, every term of $\Cov(e_3,e_1)=\sum_nW_{33}^nW_{31}W_{11}^{n+1}\Var(e_1)$ is nonnegative and the term $n=0$ is positive; likewise $\Sigma^c_{23}>0$. The cofactor formula for the inverse of the $3\times3$ matrix $\Sigma^c$ gives
\[
\Pi^c_{12}=\frac{\Sigma^c_{13}\Sigma^c_{23}-\Sigma^c_{12}\Sigma^c_{33}}{\det\Sigma^c}=\frac{\Sigma^c_{13}\Sigma^c_{23}}{\det\Sigma^c}>0.
\]
\end{proof}
% src: B1-X1(a) proof, restricted to three agents per C2 S1

\begin{proof}[Proof of \cref{lem:col-kron}]
The matrix $\Pi_{\bar C\bar C}$ is positive definite with nonpositive off-diagonal entries. Write it as $\mathsf H-\mathsf N$ with $\mathsf H$ diagonal and positive and $\mathsf N\ge0$ entrywise. Since $\mathsf H^{-1/2}\Pi_{\bar C\bar C}\mathsf H^{-1/2}=I-\mathsf H^{-1/2}\mathsf N\mathsf H^{-1/2}\succ0$ and the spectral radius of the nonnegative symmetric matrix $\mathsf H^{-1/2}\mathsf N\mathsf H^{-1/2}$ is its largest eigenvalue, $\rho(\mathsf H^{-1}\mathsf N)<1$, and $\Pi_{\bar C\bar C}^{-1}=\sum_{l\ge0}(\mathsf H^{-1}\mathsf N)^l\mathsf H^{-1}\ge0$ entrywise. Since $\Pi_{C\bar C}=-\mathcal A_{C\bar C}\le0$, the correction $\mathcal A_{C\bar C}\Pi_{\bar C\bar C}^{-1}\mathcal A_{\bar CC}$ is entrywise nonnegative, which gives the couplings. For the row sums, $\Pi_{CC}\one_C=\varepsilon^\Pi_C+a^\partial_C$, $\mathcal A_{\bar CC}\one_C=a^\partial_{\bar C}$, the vector of couplings of outsiders to $C$, and $\Pi_{\bar C\bar C}\one_{\bar C}=\varepsilon^\Pi_{\bar C}+a^\partial_{\bar C}$. Hence $\Pi_{\bar C\bar C}^{-1}a^\partial_{\bar C}=\one_{\bar C}-\Pi_{\bar C\bar C}^{-1}\varepsilon^\Pi_{\bar C}$ and, with $\mathcal A_{C\bar C}\one_{\bar C}=a^\partial_C$,
\[
(\Sigma_{CC})^{-1}\one_C=\varepsilon^\Pi_C+a^\partial_C-\mathcal A_{C\bar C}(\one_{\bar C}-\Pi_{\bar C\bar C}^{-1}\varepsilon^\Pi_{\bar C})=\varepsilon^\Pi_C+\mathcal A_{C\bar C}\Pi_{\bar C\bar C}^{-1}\varepsilon^\Pi_{\bar C}.
\]
Both $\Pi_{\bar C\bar C}^{-1}\varepsilon^\Pi_{\bar C}$ and $\Pi_{\bar C\bar C}^{-1}a^\partial_{\bar C}$ are nonnegative and add to $\one_{\bar C}$, so $0\le\Pi_{\bar C\bar C}^{-1}\varepsilon^\Pi_{\bar C}\le\one_{\bar C}$, which gives the bracket for the leaks. For the contrast, $\phi_C^\top(\Sigma_{CC})^{-1}\phi_C\le\phi_C^\top\Pi_{CC}\phi_C=\xout(C)$, and $(\Sigma_{CC})^{-1}-(\diag(\varepsilon^\Pi_C)+L_C)$ is a nonnegative diagonal matrix plus a Laplacian of nonnegative weights, hence positive semidefinite, so the internal rate does not decrease.
\end{proof}
% src: B1-L1 proof

The amended baseline and the results on whole populations use contrast at uniform leak in the following form.

\begin{lemma}[Contrast with uniform leak]\label{lem:col-uniform-leak}
Let $\Pi=\varepsilon^\Pi I+L^\Pi$ be diffusive with uniform leak $\varepsilon^\Pi>0$ on a finite population $V$, and let $C\subseteq V$ be a region of $n\ge2$ members.
\begin{enumerate}
\item[(a)] $r(C)=(\varepsilon^\Pi+\lambda_2(L_C))/(\varepsilon^\Pi+\cut(C)/n)$.
\item[(b)] If the couplings inside $C$ do not connect $C$, then $r(C)\le1$, with equality only if $\cut(C)=0$.
\item[(c)] The whole population has $r(V)=1+\lambda_2(L^\Pi)/\varepsilon^\Pi\ge1$.
\item[(d)] If $\cut(C)>0$, then $r(C)\le\max\{1,r_\star(C)\}$, and $r(C)<r_\star(C)$ when $r_\star(C)>1$.
\end{enumerate}
Hence the supremum \eqref{eq:col-baseline} for a finite closed comparison medium with uniform leak is attained by its whole population or by a region that its couplings connect.
\end{lemma}
% src: H6-L2
\begin{proof}
(a) With uniform leak, \eqref{eq:col-rates} gives $\xout(C)=\varepsilon^\Pi+\cut(C)/n$, and $\xin(C)$ is $\varepsilon^\Pi$ plus the smallest Rayleigh quotient of $L_C$ on $\one_C^\perp$, which is $\lambda_2(L_C)$. (b) If the couplings inside $C$ do not connect it, the indicators of a connected component of $C$ and of the rest of $C$ both lie in the kernel of $L_C$, so $\lambda_2(L_C)=0$ and $r(C)=\varepsilon^\Pi/(\varepsilon^\Pi+\cut(C)/n)$. (c) $\cut(V)=0$. (d) With $x=\lambda_2(L_C)$ and $y=\cut(C)/n>0$, $(\varepsilon^\Pi+x)/(\varepsilon^\Pi+y)$ is the convex combination of $1$ and $x/y$ with the positive weights $\varepsilon^\Pi/(\varepsilon^\Pi+y)$ and $y/(\varepsilon^\Pi+y)$. The last statement follows from (b) and (c), since a region that its couplings do not connect has contrast at most $1\le r(V)$.
\end{proof}

\begin{proof}[Derivations for \cref{rem:col-declared-geometry}]
A clique of $n$ agents with coupling $a$ has $L_C=a(nI-\one\one^\top)$, so $\lambda_2(L_C)=na$, and no coupling leaves it, so \cref{lem:col-uniform-leak}(a) gives $r=1+na/\varepsilon^\Pi$. Under the graph clause every declared edge carries $a$, so the comparison medium is the population itself, and its baseline is at least the contrast of each of its regions. Under the exchangeable clause the mean coupling over the $n(2n-1)$ pairs is $n(n-1)a/(n(2n-1))=(n-1)a/(2n-1)$. In an exchangeable population of $N$ agents with coupling $\bar a$ and leak $\varepsilon^\Pi$, a region of $n'<N$ agents has $\lambda_2=n'\bar a$ and cut $(N-n')\bar a$ per member, so its contrast $(\varepsilon^\Pi+n'\bar a)/(\varepsilon^\Pi+(N-n')\bar a)$ increases with $n'$ and is at most $(\varepsilon^\Pi+(N-1)\bar a)/(\varepsilon^\Pi+\bar a)$, which is below $1+N\bar a/\varepsilon^\Pi$, the contrast of the whole population, because $\varepsilon^\Pi(\varepsilon^\Pi+(N-1)\bar a)<(\varepsilon^\Pi+N\bar a)(\varepsilon^\Pi+\bar a)$. The baseline is therefore $1+N\bar a/\varepsilon^\Pi$, which for $N=2n$ is $1+2n(n-1)a/((2n-1)\varepsilon^\Pi)$, below $1+na/\varepsilon^\Pi$ because $2(n-1)<2n-1$. With uniform leak the admissible levels for a clique are then $(\bar r_{\mathrm{hom}},1+na/\varepsilon^\Pi]$. The numbers for the cliques joined by a matching are computed in \cref{chk:col-declared}.
\end{proof}
% src: H6-E1 (closed forms derived here), H4-P7 (monotonicity in the exchangeable population)

\begin{proof}[Proof of \cref{prop:col-whole-population}]
(a) The comparison medium is closed, so its whole population is one of its regions; its leak is uniform and its cut is zero, so \cref{lem:col-uniform-leak}(c) gives $\bar r_{\mathrm{hom}}\ge1+\lambda_2(L_{\mathrm{hom}})/\bar\varepsilon$. The population has $\cut(V)=0$, so $\xout(V)=\one^\top\Pi\one/N=\bar\varepsilon$ and $r(V)=\xin(V)/\bar\varepsilon$, with $\xin(V)$ the minimum of the Rayleigh quotient of $\Pi$ on $\one^\perp$ by \eqref{eq:col-rates}. If $V$ is a collective at level $r_0$, then $r(V)\ge r_0>\bar r_{\mathrm{hom}}$, which gives the inequality; with uniform leak $\xin(V)=\varepsilon^\Pi+\lambda_2(L^\Pi)$ and $\bar\varepsilon=\varepsilon^\Pi$.

(b) For a permutation $g$ of $V$ let $(P_gh)_i=h_{g^{-1}(i)}$ for vectors $h$, and put $\bar\Pi=|\mathsf G|^{-1}\sum_{g\in\mathsf G}P_g\Pi P_g^\top$, whose $(i,j)$ entry is the average of $\Pi_{g^{-1}(i)g^{-1}(j)}$ over $\mathsf G$. Let $i\neq j$. If $\{i,j\}$ lies in a set $\mathcal S$ of pairs on which $\mathsf G$ acts transitively, the orbit-stabilizer theorem gives every pair of $\mathcal S$ as $\{g^{-1}(i),g^{-1}(j)\}$ for the same number of elements $g$, so $\bar\Pi_{ij}$ is minus the mean coupling over $\mathcal S$, which is the coupling of the comparison medium. Otherwise $\{i,j\}$ lies in no averaged set, which happens only under the graph clause for a pair outside $E_{\mathrm d}$; $\mathsf G$ maps it to pairs outside $E_{\mathrm d}$, all of coupling zero in $\Pi$, so $\bar\Pi_{ij}=0$, as in the comparison medium. The row sums are $(\bar\Pi\one)_i=|\mathsf G|^{-1}\sum_g\varepsilon^\Pi_{g^{-1}(i)}$, because $P_g^\top\one=\one$ and $\Pi\one=\varepsilon^\Pi$; they equal $\bar\varepsilon$ when $\mathsf G$ is transitive on $V$, by the same count, and when the leak is uniform. Hence $\bar\Pi=\Pi_{\mathrm{hom}}$. The map $\Pi\mapsto\xin(V)=\min\{f^\top\Pi f:f\perp\one,\Vert f\Vert=1\}$ is a minimum of linear functions and therefore concave, and it is invariant under $\Pi\mapsto P_g\Pi P_g^\top$, because $P_g$ is orthogonal and fixes $\one$. Thus $\xin(V;\Pi_{\mathrm{hom}})\ge|\mathsf G|^{-1}\sum_g\xin(V;P_g\Pi P_g^\top)=\xin(V;\Pi)$, while $\xin(V;\Pi_{\mathrm{hom}})=\bar\varepsilon+\lambda_2(L_{\mathrm{hom}})$, and (a) excludes $V$. The symmetric group acts transitively on agents and on pairs. On a torus the translations map the torus to itself and act transitively on the sites and on each displacement class, every pair of the class of $x$ being $\{y,y+x\}$ for a site $y$. For the graph clause the hypothesis is the transitivity of the automorphisms of the declared graph on its vertices and on its edges.

(c) Since $\lambda_2(L^\Pi)>\lambda_2(L_{\mathrm{hom}})\ge0$, the couplings of $\Pi$ connect $V$. Every pair of positive coupling in $\Pi$ has positive coupling in the comparison medium, whose coupling on a displacement class, on $E_{\mathrm d}$ or on all pairs is a mean of nonnegative numbers one of which is positive; so the comparison medium is connected, every proper region of it has a positive cut, and $R^\star_{\mathrm{hom}}$, a maximum over finitely many regions, is finite. Let $C\subsetneq V$ with $|C|\ge2$ be a region of the comparison medium. If its couplings do not connect it, its contrast is at most $1$ by \cref{lem:col-uniform-leak}(b), below $r(V)=1+\lambda_2(L^\Pi)/\varepsilon^\Pi$. If they connect it and $r_\star(C)\le1$, its contrast is at most $1$ by \cref{lem:col-uniform-leak}(d); if $r_\star(C)>1$, its contrast is below $r_\star(C)\le R^\star_{\mathrm{hom}}$, and $R^\star_{\mathrm{hom}}\le r(V)$ when $\varepsilon^\Pi\le\lambda_2(L^\Pi)/(R^\star_{\mathrm{hom}}-1)$. The whole comparison medium has contrast $1+\lambda_2(L_{\mathrm{hom}})/\varepsilon^\Pi<r(V)$. The baseline, a maximum over finitely many regions, is therefore below $r(V)$, and since $\varepsilon^{\max}_V/\bar\varepsilon_V=1\le\bar r_{\mathrm{hom}}$, every level in $(\bar r_{\mathrm{hom}},r(V)]$ is admissible and $V$ is a collective at it. With the couplings fixed, (a) gives necessity at every leak and the bound just proved gives sufficiency at small leaks. At a given leak, $V$ is a collective at some admissible level exactly when $r(V)>\bar r_{\mathrm{hom}}$, and under the hypothesis $r(V)$ exceeds the contrast of the whole comparison medium, so only its proper regions remain to be compared.

(d) On vectors $(h_1,h_2,-h_2,-h_1)$ the Laplacian of the path acts as $\bigl(\begin{smallmatrix}0.9&-0.9\\-0.9&3.3\end{smallmatrix}\bigr)$, whose eigenvalues are $2.1\pm1.5$, and $(1,-1,-1,1)$ is an eigenvector with eigenvalue $1.8$, so $\lambda_2(L^\Pi)=0.6$. The comparison path has coupling $1$ between neighbors, the mean over the three pairs at displacement one, and no other coupling, so $\lambda_2(L_{\mathrm{hom}})=2-2\cos(\pi/4)$. The regions of the comparison path that its couplings connect are the two end pairs, with $\lambda_2=2$ and $\cut/n=\frac12$, the middle pair, with $2$ and $1$, the two triples, with $1$ and $\frac13$, and the whole path; the other regions have contrast at most one. Clearing denominators, $1+0.6/\varepsilon^\Pi$ exceeds the contrast of an end pair exactly when $\varepsilon^\Pi<\frac13$, of the middle pair when $\varepsilon^\Pi<\frac32$, of a triple when $\varepsilon^\Pi<3$, and that of the whole path at every leak.
\end{proof}
% src: H6-P1 proof (group average and concavity, sufficiency), H6-L2; the case analysis of (d) is written out here

\begin{proof}[Proof of \cref{thm:col-lattice-contrast}]
With uniform leak, $\xin(C)=\varepsilon^\Pi+\lambda_2(L_C)$ and $\xout(C)=\varepsilon^\Pi+\cut(C)/n$, and $(\varepsilon^\Pi+x)/(\varepsilon^\Pi+y)\le\max\{1,x/y\}$ for $x,y\ge0$, which gives $r\le\max\{1,r_\star\}$. Every line parallel to an axis that meets $C$ meets it in a finite nonempty set and so contains at least two lattice edges with exactly one endpoint in $C$, each of coupling at least $a_{\mathrm{nn}}$, and edges on different lines are different. With $\operatorname{pr}_d(C)$ the projection of $C$ along axis $d$, the number of such lines is $|\operatorname{pr}_d(C)|$, so $\cut(C)\ge2a_{\mathrm{nn}}\sum_d|\operatorname{pr}_d(C)|\ge2Da_{\mathrm{nn}}n^{(D-1)/D}$ by the inequality of arithmetic and geometric means and the discrete Loomis--Whitney inequality $n^{D-1}\le\prod_d|\operatorname{pr}_d(C)|$ \citep{loomis1949inequality}. Gershgorin's theorem gives $\lambda_2(L_C)\le2\max_i\sum_{j\in C}a^\Pi_{ij}\le2\Vert a^\Pi\Vert_1$, which yields the first term of \eqref{eq:col-lattice}. For the second, \cref{thm:col-contrast-size} gives $\lambda_2(L_C)\le\varsigma^2/(2\max_d\Var_d(C))$. A sup-norm ball of radius $h_n=((n/2)^{1/D}-1)/2$ around the centroid $\bar x$ of $C$ contains at most $(2h_n+1)^D=n/2$ lattice sites, so at least $n/2$ members lie farther than $h_n$ from $\bar x$, $\sum_d\Var_d(C)\ge h_n^2/2$ and $\max_d\Var_d(C)\ge((n/2)^{1/D}-1)^2/(8D)$. Dividing by $\cut(C)/n\ge2Da_{\mathrm{nn}}n^{-1/D}$ gives the second term. For $n\le2^{D+1}$ the first term is at most $2^{(D+1)/D}\Vert a^\Pi\Vert_1/(Da_{\mathrm{nn}})$; for $n>2^{D+1}$, $(n/2)^{1/D}-1\ge\frac12(n/2)^{1/D}$ and the second term is at most $2^{3+2/D}\varsigma^2n^{-1/D}/a_{\mathrm{nn}}$. Both bounds are uniform in $\varepsilon^\Pi$.
\end{proof}
% src: B1-P2 proof, L2-T3 proof (explicit uniform constant added here)

\begin{proof}[Proof of \cref{prop:col-block-media}]
(a) The Laplacian of an $\ell_b\times\ell_b$ grid is the Kronecker sum of two path Laplacians with eigenvalues $2-2\cos(\pi j/\ell_b)$, $j=0,\dots,\ell_b-1$, so $\lambda_2(L_C)=a_{\mathrm{in}}(2-2\cos(\pi/\ell_b))$. A block has $4\ell_b$ edges to other blocks, one for each site on each side, so $\cut(C)/n=4\ell_ba_{\mathrm{out}}/\ell_b^2$. (b) A vector $f=\sum_jc_j\one_{C_j}$ in the span of the block indicators has
\[
f^\top L^\Pi f=\sum_{\{i,k\}\text{ across blocks}}a^\Pi_{ik}(c_{j(i)}-c_{j(k)})^2\le2\sum_jc_j^2\cut(C_j)\le2\max_j\frac{\cut(C_j)}{|C_j|}\Vert f\Vert^2,
\]
and Courant--Fischer on this $M$-dimensional space gives the upper bound. For the lower bound, $L^\Pi=\bigoplus_jL_{C_j}+L_\times$ with $L_\times\succeq0$ the Laplacian of the couplings between blocks; the direct sum has exactly $M$ zero eigenvalues, because the blocks are connected, and next eigenvalue $\min_j\lambda_2(L_{C_j})$, and adding $L_\times$ lowers no eigenvalue.
\end{proof}
% src: L2-P9 proof

\begin{proof}[Proof of \cref{thm:col-contrast-size}]
The function $f=x_d-\bar x_{C,d}$ on $C$, zero elsewhere, is supported on $C$ and orthogonal to $\one_C$. By Courant--Fischer $\xin(C)\le f^\top(\diag(\varepsilon^\Pi)+L_C)f/f^\top f\le\varepsilon^{\max}_C+f^\top L_Cf/f^\top f$ and $\lambda_2(L_C)\le f^\top L_Cf/f^\top f$. Now
\[
f^\top L_Cf=\frac12\sum_{i,j\in C}a^\Pi_{ij}(x_{i,d}-x_{j,d})^2\le\frac12\sum_{i\in C}\varsigma_d^2=\frac{n\varsigma_d^2}{2},\qquad f^\top f=n\Var_d(C),
\]
which gives the first inequality of \eqref{eq:col-size}; dividing by $\xout(C)\ge\bar\varepsilon_C$ gives the second. For the metric form take $f=\operatorname{dist}(\cdot,j_0)$ minus its mean over $C$, for which $|f_i-f_j|\le\operatorname{dist}(i,j)$. For sharpness, an isolated box has $\xout=\varepsilon^\Pi$ and $r-1=\lambda_2(L_C)/\varepsilon^\Pi$ exactly. For a box of side $\ell_b$ with coupling $a^\Pi$, $\lambda_2(L_C)=2a^\Pi(1-\cos(\pi/\ell_b))$, $\Var_d(C)=(\ell_b^2-1)/12$ and $\xi_{2,d}^2=a^\Pi/\varepsilon^\Pi$, so $\Var_d(C)(r-1)/\xi_{2,d}^2=(\ell_b^2-1)(1-\cos(\pi/\ell_b))/6\to\pi^2/12$.
\end{proof}
% src: B1-T2 proof

\begin{proof}[Proof of \cref{cor:col-lstar}]
Let $r(C)\ge r_0$. The second inequality of \eqref{eq:col-size} gives $r_0\bar\varepsilon_C-\varepsilon^{\max}_C\le\varsigma_d^2/(2\Var_d(C))$, and the left side is positive by \eqref{eq:col-collective}, so $\Var_d(C)\le\varsigma_d^2/(2(r_0\bar\varepsilon_C-\varepsilon^{\max}_C))$. With uniform leak the right side is $\xi_{2,d}^2/(r_0-1)$, and the metric form is the same computation. With heterogeneous leaks $r_0\bar\varepsilon_C-\varepsilon^{\max}_C\ge r_0\varepsilon^\Pi_{\min}-\varepsilon^\Pi_{\max}$. Taking suprema over collectives gives the corollary.
\end{proof}
% src: B1-C1 proof

The proof of the bridge theorem uses two more facts.

\begin{lemma}[Kato--Temple]\label{lem:col-kato-temple}
Let $\mathsf M$ be a real symmetric operator on a finite-dimensional space with smallest eigenvalues $\lambda_1(\mathsf M)\le\lambda_2(\mathsf M)$, let $u$ be a unit vector with Rayleigh quotient $\langle\mathsf M\rangle_u=u^\top\mathsf Mu<\lambda_2(\mathsf M)$, and put $\mathsf r_u=\Vert(\mathsf M-\langle\mathsf M\rangle_u)u\Vert$. Then $0\le\langle\mathsf M\rangle_u-\lambda_1(\mathsf M)\le\mathsf r_u^2/(\lambda_2(\mathsf M)-\langle\mathsf M\rangle_u)$.
\end{lemma}
% src: L2-L3(b)
\begin{proof}
Write $\lambda_1,\lambda_2$ for $\lambda_1(\mathsf M),\lambda_2(\mathsf M)$. The Rayleigh principle gives $\lambda_1\le\langle\mathsf M\rangle_u$. The operator $(\mathsf M-\lambda_1)(\mathsf M-\lambda_2)$ is positive semidefinite, its eigenvalues being $(\lambda_k(\mathsf M)-\lambda_1)(\lambda_k(\mathsf M)-\lambda_2)\ge0$, and since $(\mathsf M-\langle\mathsf M\rangle_u)u\perp u$, $\Vert\mathsf Mu\Vert^2=\mathsf r_u^2+\langle\mathsf M\rangle_u^2$. Hence $0\le u^\top(\mathsf M-\lambda_1)(\mathsf M-\lambda_2)u=\mathsf r_u^2-(\langle\mathsf M\rangle_u-\lambda_1)(\lambda_2-\langle\mathsf M\rangle_u)$.
\end{proof}
This is the inequality of \citet{temple1928theory} and \citet{kato1949upper}.

\begin{lemma}[Effective resistance]\label{lem:col-resistance}
For a connected population with Laplacian $L^\Pi$ and agents $i\neq j$, let $R_{\mathrm{eff}}(i,j)=(\one_{\{i\}}-\one_{\{j\}})^\top(L^\Pi)^+(\one_{\{i\}}-\one_{\{j\}})$. Then $1/R_{\mathrm{eff}}(i,j)=\min\{f^\top L^\Pi f:f_i=1,f_j=0\}$. Consequently $R_{\mathrm{eff}}(i,j)\ge1/\cut(C)$ when $i\in C$ and $j\notin C$, and $R_{\mathrm{eff}}(i,j)\le2/\lambda_2(L_C)$ when $i,j\in C$ and $C$ is connected by its internal couplings.
\end{lemma}
% src: B1-T3(iii)(a) proof (Rayleigh monotonicity and shorting, proved here from the Dirichlet principle)
\begin{proof}
Since $L^\Pi\one=0$, the constraint can be replaced by $f_i-f_j=1$. With $b=\one_{\{i\}}-\one_{\{j\}}$, which lies in the range of $L^\Pi$, a minimizer satisfies $L^\Pi f=\mu b$, so $f=\mu(L^\Pi)^+b$ up to a constant, the constraint gives $\mu=1/R_{\mathrm{eff}}(i,j)$, and the minimum is $\mu^2b^\top(L^\Pi)^+b=1/R_{\mathrm{eff}}(i,j)$. The indicator $f=\one_C$ is admissible when $i\in C$, $j\notin C$, and has $f^\top L^\Pi f=\cut(C)$, which gives the first consequence. For $i,j\in C$, $f^\top L^\Pi f\ge f_C^\top L_Cf_C$ for every $f$, so the minimum over the population is at least the minimum for the couplings inside $C$, whose value is $1/(b_C^\top L_C^+b_C)$ with $b_C$ the restriction of $b$; since $b_C\perp\one_C$ and $\Vert b_C\Vert^2=2$, $b_C^\top L_C^+b_C\le2/\lambda_2(L_C)$.
\end{proof}

\begin{proof}[Proof of \cref{thm:col-bridge}]
(i) With uniform leak $\xin(C_j)=\varepsilon^\Pi+g_j$ and $\xout(C_j)=\varepsilon^\Pi+\bar a^\partial_j$. When two blocks form $V$, $\cut(V)=0$, so $\xout(V)=\varepsilon^\Pi$ and $\xin(V)=\varepsilon^\Pi+\lambda_2(L^\Pi)$. The unit vector $u_-=(\phi_{C_1}-\phi_{C_2})/\sqrt2$ is orthogonal to $\one$, takes the values $\pm(2n)^{-1/2}$, and has $u_-^\top L^\Pi u_-=\cut(C_1)\cdot4/(2n)=2\bar a^\partial$, so $\lambda_2(L^\Pi)\le2\bar a^\partial$ by Courant--Fischer on $\one^\perp$. The vector $(L^\Pi-2\bar a^\partial)u_-$ has entries $\pm2(a^\partial_i-\bar a^\partial)/\sqrt{2n}$, with sign $+$ on $C_1$ and $-$ on $C_2$; its block sums vanish, so it is orthogonal to the span of $\Phi$, and its squared norm is $\frac2n\sum_i(a^\partial_i-\bar a^\partial)^2$. Since $L^\Pi\Phi v_+=L^\Pi\one/\sqrt N=0$ for $v_+=(1,1)/\sqrt2$ and $L^\Pi u_-$ equals $2\bar a^\partial u_-$ plus this vector, $\sigma^2_{\mathrm{cut}}=\Vert\Phi_\perp^\top L^\Pi\Phi\Vert_F^2$ equals that squared norm. On $\one^\perp$ the second eigenvalue of $L^\Pi$ is $\lambda_3(L^\Pi)\ge g_{\min}$ by \cref{prop:col-block-media}(b) with $\varepsilon^\Pi=0$. If $2\bar a^\partial<g_{\min}$, \cref{lem:col-kato-temple} on $\one^\perp$ with $u=u_-$ gives $2\bar a^\partial-\lambda_2(L^\Pi)\le\sigma^2_{\mathrm{cut}}/(g_{\min}-2\bar a^\partial)$, which is \eqref{eq:col-union}.

(ii) In the orthonormal basis $[\Phi,\Phi_\perp]$ the block formula for the inverse gives $\Phi^\top\Pi^{-1}\Phi=(\Pi_\Phi-\mathcal D)^{-1}$, and $0\preceq\mathcal D\preceq\Vert\Phi_\perp^\top\Pi\Phi\Vert^2\lambda_{\min}(\Phi_\perp^\top\Pi\Phi_\perp)^{-1}I$ with $\Vert\cdot\Vert\le\Vert\cdot\Vert_F$; the squared Frobenius norm is the sum of the squared norms of the columns $\Phi_\perp^\top\Pi\phi_{C_j}$, one per block, and for $M$ identical blocks in identical surroundings these are equal, so that $\sigma^2_{\mathrm{cut}}$ is $M$ times one of them. $\mathcal D=0$ exactly when $\Phi_\perp^\top\Pi\Phi=0$, that is when the span of $\Phi$ is invariant under $\Pi$; with uniform leak this means that $L^\Pi\one_{C_j}$ is constant on every block for every $j$, and since $(L^\Pi\one_{C_j})_i=-\sum_{k\in C_j}a^\Pi_{ik}$ for $i\notin C_j$ and $\sum_{l\neq j}\sum_{k\in C_l}a^\Pi_{ik}$ for $i\in C_j$, this is the equitable condition. For two equal blocks, $\Pi_\Phi=\begin{pmatrix}\varepsilon^\Pi+\bar a^\partial&-\bar a^\partial\\-\bar a^\partial&\varepsilon^\Pi+\bar a^\partial\end{pmatrix}$ has eigenvectors $v_\pm=(1,\pm1)/\sqrt2$ with eigenvalues $\varepsilon^\Pi$ and $\varepsilon^\Pi+2\bar a^\partial$. Because $\Phi v_+=\one/\sqrt N$ and $\Pi\one=\varepsilon^\Pi\one$, $\Phi_\perp^\top\Pi\Phi v_+=0$, so $\Phi_\perp^\top\Pi\Phi$ has rank at most one, its spectral and Frobenius norms coincide, and $\mathcal D=\mathcal D_{--}v_-v_-^\top$ with $0\le\mathcal D_{--}\le\sigma^2_{\mathrm{cut}}/\lambda_{\min}(\Phi_\perp^\top\Pi\Phi_\perp)$. Here $\Phi_\perp^\top\Pi\Phi_\perp\succeq\varepsilon^\Pi I+\Phi_\perp^\top(L_{C_1}\oplus L_{C_2})\Phi_\perp\succeq(\varepsilon^\Pi+g_{\min})I$, since the columns of $\Phi_\perp$ are orthogonal to the kernel of $L_{C_1}\oplus L_{C_2}$, so $\mathcal D_{--}\le\delta_{\mathrm{corr}}$. The matrix $\Phi^\top\Sigma\Phi$ then has eigenvalues $1/\varepsilon^\Pi$ on $v_+$ and $1/(\varepsilon^\Pi+2\bar a^\partial-\mathcal D_{--})$ on $v_-$, and the correlation of $\phi_{C_1}^\top e$ and $\phi_{C_2}^\top e$, which is that of the block averages, is
\[
\Corr(\bar e_1,\bar e_2)=\frac{2\bar a^\partial-\mathcal D_{--}}{2\varepsilon^\Pi+2\bar a^\partial-\mathcal D_{--}},
\]
decreasing in $\mathcal D_{--}$, which gives \eqref{eq:col-corr}.

(iii)(a) For the closed connected population, $\Sigma=(\varepsilon^\Pi I+L^\Pi)^{-1}=\one\one^\top/(N\varepsilon^\Pi)+(L^\Pi)^++O(\varepsilon^\Pi)$ by the spectral decomposition. Write $\Sigma_{ij}=\alpha_0+\alpha_{ij}$ with $\alpha_0=1/(N\varepsilon^\Pi)$, $\alpha_{ij}=O(1)$ and $\alpha_i=\alpha_{ii}$. Expanding in $1/\alpha_0$, $\Corr(e_i,e_j)=1-(\alpha_i+\alpha_j-2\alpha_{ij})/(2\alpha_0)+O(\alpha_0^{-2})$, and $\alpha_i+\alpha_j-2\alpha_{ij}=(L^\Pi)^+_{ii}+(L^\Pi)^+_{jj}-2(L^\Pi)^+_{ij}+O(\varepsilon^\Pi)=R_{\mathrm{eff}}(i,j)+O(\varepsilon^\Pi)$. Hence $\varpi_{ij}=\frac{N\varepsilon^\Pi}2R_{\mathrm{eff}}(i,j)+O((\varepsilon^\Pi)^2)$ uniformly over the finitely many pairs. The single-linkage ultrametric, $u_{\mathrm{in}}$ and $\varpi_{\mathrm{out}}$ change by no more than the largest change of the weights, and they scale with the weights, so $\mathcal U(C)=\frac{N\varepsilon^\Pi}2\mathcal U^R(C)+O((\varepsilon^\Pi)^2)$. By \cref{lem:col-resistance}, the resistance version of $\varpi_{\mathrm{out}}(C)$ is at least $1/\cut(C)$, and the resistance version of $u_{\mathrm{in}}(C)$ is at most $\max_{i,j\in C}R_{\mathrm{eff}}(i,j)\le2/\lambda_2(L_C)=2|C|/(r_\star(C)\cut(C))$, because a single-linkage height never exceeds the direct weight.

(iii)(b) Write $L^\Pi=\operatorname{Deg}-\mathcal A$ with $\operatorname{Deg}$ the diagonal of coupling sums and $\mathcal A$ the coupling matrix. Since $(\varepsilon^\Pi I+\operatorname{Deg})^{-1}\mathcal A$ is nonnegative with row sums below one, $\Sigma=\sum_{l\ge0}((\varepsilon^\Pi I+\operatorname{Deg})^{-1}\mathcal A)^l(\varepsilon^\Pi I+\operatorname{Deg})^{-1}$ with nonnegative terms. As $\varepsilon^\Pi\to\infty$, $\Sigma_{ii}=(\varepsilon^\Pi)^{-1}+O((\varepsilon^\Pi)^{-2})$, a coupled pair has $\Corr(e_i,e_j)=a^\Pi_{ij}/\varepsilon^\Pi+O((\varepsilon^\Pi)^{-2})$ and $\varpi_{ij}=\ln(\varepsilon^\Pi/a^\Pi_{ij})+O(1/\varepsilon^\Pi)$, and a pair at distance $l\ge2$ in the coupling graph has $\Corr=O((\varepsilon^\Pi)^{-l})$ and $\varpi_{ij}\ge l\ln\varepsilon^\Pi-O(1)$. For large $\varepsilon^\Pi$ every chain through a non-coupled pair is therefore worse than the best chain of coupled pairs, and the single-linkage height between $i$ and $j$ is $\ln\varepsilon^\Pi-\ln a^{\mathrm{bn}}(i,j)+O(1/\varepsilon^\Pi)$, with $a^{\mathrm{bn}}(i,j)$ the largest, over paths from $i$ to $j$, of the smallest coupling on the path. Inside $C$ every pair is joined in a spanning tree of maximal weight through couplings at least $a^{\mathrm{bn}}(C)$; removing the tree edge of coupling $a^{\mathrm{bn}}(C)$ splits $C$ into two parts that no edge inside $C$ of larger coupling joins, by the cut property of such trees, and any path that leaves $C$ crosses a boundary coupling at most $a^\partial_{\max}(C)<a^{\mathrm{bn}}(C)$. Hence $u_{\mathrm{in}}(C)=\ln\varepsilon^\Pi-\ln a^{\mathrm{bn}}(C)+O(1/\varepsilon^\Pi)$. The smallest direct weight across the boundary is at the boundary pair of largest coupling, $\varpi_{\mathrm{out}}(C)=\ln\varepsilon^\Pi-\ln a^\partial_{\max}(C)+O(1/\varepsilon^\Pi)$. Meanwhile $r(C)=(\varepsilon^\Pi+\lambda_2(L_C))/(\varepsilon^\Pi+\cut(C)/|C|)\to1$.

(iv) The statement on the blocks is $\varepsilon^\Pi+g\ge r_0(\varepsilon^\Pi+\bar a^\partial)$ rearranged. The statements on the union are the bound $r(C_1\cup C_2)\le1+2\bar a^\partial/\varepsilon^\Pi$ of (i) and, when $2\bar a^\partial<g$, the lower bound of \eqref{eq:col-union}, each solved for $\varepsilon^\Pi$. The statements on the correlation are the two sides of \eqref{eq:col-corr} solved for $\varepsilon^\Pi$ at fixed $\delta_{\mathrm{corr}}$. For the closed form put $c=2\bar a^\partial-\sigma^2_{\mathrm{cut}}/g$. If $c\le0$ the condition holds trivially. If $c>0$, then $\delta_{\mathrm{corr}}\le\sigma^2_{\mathrm{cut}}/g$ gives $2\bar a^\partial-\delta_{\mathrm{corr}}\ge c$, and since $x\mapsto x/(2\varepsilon^\Pi+x)$ increases on $x>0$, the left side of \eqref{eq:col-corr} is at least $c/(2\varepsilon^\Pi+c)$; the requirement $c/(2\varepsilon^\Pi+c)\le\theta_0$ is $\varepsilon^\Pi\ge c(1-\theta_0)/(2\theta_0)$. On the window \eqref{eq:col-window} the blocks are selected, the union is not, and the block averages are separated.
\end{proof}
% src: B1-T3 proof, L2-L3(b), L2-P9(b), C2-3

The baseline of \cref{obs:col-bridge-baseline} rests on enumeration for small regions and on the following bounds for the others.

\begin{lemma}[Bounds for regions of a homogeneous box]\label{lem:col-box-bounds}
Let a comparison medium be a box of $n_{\mathrm{row}}$ rows and $n_{\mathrm{col}}$ columns of $\Z^2$, $N=n_{\mathrm{row}}n_{\mathrm{col}}$, with coupling $a_{\mathrm h}$ between horizontal neighbors, $a_{\mathrm v}$ between vertical neighbors and no other coupling, and with Laplacian $L_{\mathrm{hom}}$. Let $S$ be a region of $n$ sites, $2\le n\le N-1$, put $k=N-n$ and $a_{\max}=\max\{a_{\mathrm h},a_{\mathrm v}\}$, and let $L_S$ be the Laplacian of the couplings inside $S$.
\begin{enumerate}
\item[(U1)] $\lambda_2(L_S)\le a_{\max}\mathsf p(n)/\mathsf J(n)$, with $\mathsf p(n)=2n-\lceil2\sqrt n\rceil$ and $\mathsf J(n)=n^2/(2\pi)-n/6$.
\item[(U2)] $\lambda_2(L_S)\le\lambda_{2+k}(L_{\mathrm{hom}})$, the $(2+k)$th smallest eigenvalue of $L_{\mathrm{hom}}$.
\item[(U3)] If $\mathsf v$ is a unit eigenvector of $L_{\mathrm{hom}}$ for $\lambda_2(L_{\mathrm{hom}})$, $\mathsf s_k$ the sum of the $k$ largest $\mathsf v_i^2$ and $\mathsf t_k$ the sum of the $k$ largest $|\mathsf v_i|$, then $\lambda_2(L_S)\le\lambda_2(L_{\mathrm{hom}})/(1-\mathsf s_k-\mathsf t_k^2/n)$ when the denominator is positive.
\item[(U4)] $\lambda_2(L_S)\le\frac n{n-1}(a_{\mathrm h}+a_{\mathrm v})$.
\item[(C)] $\cut(S)\ge\cut^{\mathrm{lb}}(n)$, the smallest of the following numbers: $\mathsf f(n)$; when $n\ge n_{\mathrm{col}}$, $a_{\mathrm v}n_{\mathrm{col}}$, plus $a_{\mathrm h}$ if $n_{\mathrm{col}}$ does not divide $n$; when $n\ge n_{\mathrm{row}}$, $a_{\mathrm h}n_{\mathrm{row}}$, plus $a_{\mathrm v}$ if $n_{\mathrm{row}}$ does not divide $n$; and when $n\ge n_{\mathrm{row}}+n_{\mathrm{col}}-1$, $\mathsf f(N-n)$. Here $\mathsf f(m)$ is the minimum of $a_{\mathrm h}\mathsf y+a_{\mathrm v}\mathsf x$ over integers $1\le\mathsf y\le n_{\mathrm{row}}$ and $1\le\mathsf x\le n_{\mathrm{col}}$ with $\mathsf x\mathsf y\ge m$, $\mathsf y(n_{\mathrm{col}}-1)\ge m$ and $\mathsf x(n_{\mathrm{row}}-1)\ge m$, and $+\infty$ if there are none.
\end{enumerate}
Hence at uniform leak $\varepsilon^\Pi$ every region of $n$ sites has $r(S)\le(\varepsilon^\Pi+\lambda^{\mathrm{ub}}(n))/(\varepsilon^\Pi+\cut^{\mathrm{lb}}(n)/n)$, with $\lambda^{\mathrm{ub}}(n)$ the smallest of the bounds (U1) to (U4).
\end{lemma}
% src: H6-L3 (symbols renamed to avoid reserved letters)
\begin{proof}
Give the sites coordinates $z_i=(x_i,y_i)$, column and row. (U1) The vector $f$ on $S$ with $f_i=x_i-\bar x_S$, $\bar x_S$ the mean column of $S$, is orthogonal to $\one_S$, $f^\top f=\sum_{i\in S}(x_i-\bar x_S)^2$, and $f^\top L_Sf=a_{\mathrm h}E_{\mathrm h}$, with $E_{\mathrm h}$ the number of horizontal neighbor pairs inside $S$, since such pairs differ by one in column and vertical pairs by zero. The row coordinate gives the quotient of $a_{\mathrm v}E_{\mathrm v}$ by $\sum_{i\in S}(y_i-\bar y_S)^2$ in the same way. By the mediant inequality $\min\{p_1/q_1,p_2/q_2\}\le(p_1+p_2)/(q_1+q_2)$, where a vanishing denominator comes with a vanishing numerator and the other quotient is used, $\lambda_2(L_S)\le a_{\max}(E_{\mathrm h}+E_{\mathrm v})/\mathsf J_S$ with $\mathsf J_S=\sum_{i\in S}|z_i-\bar z_S|^2$. If $S$ meets $\mathsf y$ rows and $\mathsf x$ columns, a row holding $m'$ sites of $S$ holds at most $m'-1$ horizontal pairs inside $S$, so $E_{\mathrm h}\le n-\mathsf y$, and likewise $E_{\mathrm v}\le n-\mathsf x$; since $\mathsf x\mathsf y\ge n$, $\mathsf x+\mathsf y\ge2\sqrt n$, and both are integers, so $E_{\mathrm h}+E_{\mathrm v}\le\mathsf p(n)$. For the inertia, put the unit square $Q_z$ centered at each site $z\in S$. These squares have disjoint interiors, $\int_{Q_z}|w-c|^2\mathrm dw=|z-c|^2+\frac16$, and their union $A_S$ has area $n$. If $B$ is the disk of area $n$ centered at $c$, with radius $\varrho$, then $|w-c|\ge\varrho$ on $A_S\setminus B$, $|w-c|\le\varrho$ on $B\setminus A_S$, and the two sets have equal areas, so $\int_{A_S}|w-c|^2\mathrm dw\ge\int_B|w-c|^2\mathrm dw=n^2/(2\pi)$. With $c=\bar z_S$ this is $\mathsf J_S+n/6\ge n^2/(2\pi)$, which gives (U1).

(U2) The principal submatrix of $L_{\mathrm{hom}}$ on $S$ is $L_S$ plus the diagonal matrix of the members' cut couplings, so $\lambda_2(L_S)$ is at most its second eigenvalue, and Cauchy interlacing bounds that by $\lambda_{2+N-n}(L_{\mathrm{hom}})$.

(U3) Put $f=\mathsf v|_S-\bar{\mathsf v}_S\one_S$, with $\bar{\mathsf v}_S$ the mean of $\mathsf v$ on $S$. Then $f\perp\one_S$, and $f^\top L_Sf$ is the sum of $a_{ij}(\mathsf v_i-\mathsf v_j)^2$ over the neighbor pairs inside $S$, at most the same sum over all pairs, which is $\lambda_2(L_{\mathrm{hom}})$. Since $\mathsf v\perp\one$, $|\sum_{i\in S}\mathsf v_i|=|\sum_{i\notin S}\mathsf v_i|\le\mathsf t_k$, and $\sum_{i\in S}\mathsf v_i^2\ge1-\mathsf s_k$, so $f^\top f=\sum_{i\in S}\mathsf v_i^2-(\sum_{i\in S}\mathsf v_i)^2/n\ge1-\mathsf s_k-\mathsf t_k^2/n$.

(U4) For $i\in S$ the vector $\one_{\{i\}}-\one_S/n$ is orthogonal to $\one_S$, has squared norm $(n-1)/n$, and has quadratic form $(L_S)_{ii}$, the coupling of $i$ inside $S$. Take for $i$ the site of $S$ with the smallest column in the lowest row that meets $S$: it has no neighbor in $S$ to its left or below, so its coupling inside $S$ is at most $a_{\mathrm h}+a_{\mathrm v}$.

(C) A row that meets $S$ without being contained in it contains a horizontal pair with exactly one site in $S$, and such pairs in different rows are different; the same holds for columns and vertical pairs. If $S$ contains no full row and no full column, each of the $\mathsf y$ rows and $\mathsf x$ columns that it meets contributes to the cut, a row holds at most $n_{\mathrm{col}}-1$ of its sites and a column at most $n_{\mathrm{row}}-1$, and $\mathsf x\mathsf y\ge n$, so $\cut(S)\ge a_{\mathrm h}\mathsf y+a_{\mathrm v}\mathsf x\ge\mathsf f(n)$. If $S$ contains a full row and no full column, then $n\ge n_{\mathrm{col}}$ and all $n_{\mathrm{col}}$ columns meet $S$ without being contained in it, so $\cut(S)\ge a_{\mathrm v}n_{\mathrm{col}}$, plus $a_{\mathrm h}$ when $n_{\mathrm{col}}$ does not divide $n$, since then some row that meets $S$ is not contained in it; the case of a full column and no full row is symmetric. If $S$ contains a full row and a full column, then $n\ge n_{\mathrm{row}}+n_{\mathrm{col}}-1$ and $S$ meets every row and every column, so every row or column that meets the complement $\bar S$ of $S$ also meets $S$; then $\bar S$, with $N-n$ sites, is in the first case, and $\cut(S)=\cut(\bar S)\ge\mathsf f(N-n)$. The last statement holds because at uniform leak $r(S)$ increases with $\lambda_2(L_S)$ and decreases with $\cut(S)$, by \cref{lem:col-uniform-leak}(a).
\end{proof}
% src: H6-L3 proof; the bathtub step is proved here directly

\begin{proof}[Proof of \cref{prop:col-coherence-average}]
(a) By the Cauchy--Schwarz inequality, $1=(\phi_C^\top\phi_C)^2\le(\phi_C^\top\Sigma\phi_C)\allowbreak(\phi_C^\top\Pi\phi_C)$. By \cref{lem:col-kron}, $(\Sigma_{CC})^{-1}\succeq\varepsilon^\Pi I+L_C$, so $\Sigma_{CC}\preceq(\varepsilon^\Pi I+L_C)^{-1}$, whose restriction to the orthogonal complement of $\one_C$ has norm $1/(\varepsilon^\Pi+\lambda_2(L_C))=1/\xin(C)$. (b) $\sum_{i\in C}\sigma_i\Corr(e_i,\bar e_C)=\sum_i\Cov(e_i,\bar e_C)/\operatorname{sd}(\bar e_C)=n\operatorname{sd}(\bar e_C)$ and $\sum_i\sigma_i\le(n\tr\Sigma_{CC})^{1/2}$, so the left side of \eqref{eq:col-coherence} is at least $(\mathsf V/\tr\Sigma_{CC})^{1/2}$ with $\mathsf V=\Var(\phi_C^\top e)=n\Var(\bar e_C)$. Computing the trace in an orthonormal basis that contains $\phi_C$, part (a) gives $\tr\Sigma_{CC}\le\mathsf V+(n-1)/\xin(C)$, and $\mathsf V/(\mathsf V+(n-1)/\xin)$ increases in $\mathsf V\ge1/\xout$, which gives $r/(r+n-1)$. (c) For an isolated region $\Sigma=(\varepsilon^\Pi I+L_C)^{-1}$, so $\Cov(e_i,\bar e_C)=\Var(\bar e_C)=1/(n\varepsilon^\Pi)$ and $\Sigma_{ii}=1/(n\varepsilon^\Pi)+\sum_{k\ge2}v_k(i)^2/(\varepsilon^\Pi+\lambda_k(L_C))$ for orthonormal eigenvectors $v_k$. Hence $\Corr(e_i,\bar e_C)^{-2}=n\varepsilon^\Pi\Sigma_{ii}$, and averaging with $\sum_iv_k(i)^2=1$ gives the identity. For an interval with coupling $a^\Pi$, the Dirichlet principle of \cref{lem:col-resistance} gives $R_{\mathrm{eff}}(i,j)=|i-j|/a^\Pi$, the minimizer being linear between $i$ and $j$. Since $R_{\mathrm{eff}}(i,j)=L^+_{ii}+L^+_{jj}-2L^+_{ij}$ and $L^+\one=0$, $\sum_{i<j}R_{\mathrm{eff}}(i,j)=n\tr L_C^+$, so $\sum_{k\ge2}1/\lambda_k=\tr L_C^+=n^{-1}\sum_{i<j}(j-i)/a^\Pi=(n^2-1)/(6a^\Pi)$. At contrast $r$, $\varepsilon^\Pi=\lambda_2(L_C)/(r-1)$ with $\lambda_2(L_C)=2a^\Pi(1-\cos(\pi/n))\le a^\Pi\pi^2/n^2$, so $\varepsilon^\Pi T_C\le\varepsilon^\Pi\tr L_C^+\le\pi^2(n^2-1)/(6n^2(r-1))$.
\end{proof}
% src: B1-P4 proof, B1-P5 proof (Kirchhoff-index identity rederived here)

\begin{proof}[Proof of \cref{lem:col-coarse-medium}]
(a) For $\alpha\neq\beta$ the blocks are disjoint, so every term of $\Pi^{\mathcal P}_{\alpha\beta}$ is an off-diagonal entry of $\Pi$ and $\Pi^{\mathcal P}_{\alpha\beta}=-\cut(C_\alpha,C_\beta)\le0$; the row sums are $\sum_\beta\Pi^{\mathcal P}_{\alpha\beta}=\sum_{i\in C_\alpha}(\Pi\one)_i=\sum_{i\in C_\alpha}\varepsilon^\Pi_i\ge0$. The matrix $\Pi^{\mathcal P}$ is symmetric, and because the blocks are disjoint and nonempty, $\mathsf P^\top$ has full column rank, so $h\neq0$ gives $h^\top\Pi^{\mathcal P}h=(\mathsf P^\top h)^\top\Pi(\mathsf P^\top h)>0$. (b) In the orthonormal basis $[\Phi,\Phi_\perp]$ write $e=\Phi e_\Phi+\Phi_\perp e_\perp$, with $e_\Phi=\Phi^\top e$ and $e_\perp=\Phi_\perp^\top e$. The pair $(e_\Phi,e_\perp)$ is Gaussian with precision $\begin{pmatrix}\Phi^\top\Pi\Phi&\Phi^\top\Pi\Phi_\perp\\\Phi_\perp^\top\Pi\Phi&\Phi_\perp^\top\Pi\Phi_\perp\end{pmatrix}$. Given $e_\perp=0$, $e_\Phi$ is centered with precision $\Phi^\top\Pi\Phi$; without conditioning its precision is the Schur complement $\Phi^\top\Pi\Phi-\mathcal D$, with $\mathcal D\succeq0$ because $\Phi_\perp^\top\Pi\Phi_\perp\succ0$. Since $\mathsf P=\diag(n_\alpha)^{1/2}\Phi^\top$, the block averages are $\bar e=\diag(n_\alpha)^{-1/2}e_\Phi$, so their precision is $\diag(n_\alpha)^{1/2}\Phi^\top\Pi\Phi\diag(n_\alpha)^{1/2}=\Pi^{\mathcal P}$ given $e_\perp=0$ and $\Pi^{\mathcal P}-\diag(n_\alpha)^{1/2}\mathcal D\diag(n_\alpha)^{1/2}$ without conditioning. The correction vanishes, $\mathcal D=0$, exactly when $\Phi_\perp^\top\Pi\Phi=0$, that is when $\Pi$ maps the span of $\Phi$, which is the span of the block indicators, into itself; with uniform leak this is the equitable condition, as in the proof of \cref{thm:col-bridge}(ii). (c) Every block indicator of $\mathcal P_{l+1}$ is the sum of the block indicators of $\mathcal P_l$ that it contains, which is $\mathsf P_{l+1}=\mathsf Q_l\mathsf P_l$, and then $\mathsf P_{l+1}\Pi\mathsf P_{l+1}^\top=\mathsf Q_l\mathsf P_l\Pi\mathsf P_l^\top\mathsf Q_l^\top$. (d) With $n_\alpha=s$, $\Phi=s^{-1/2}\mathsf P^\top$, so $\Pi^{\mathcal P}=s\Phi^\top\Pi\Phi$, and contrast and the ratio of leaks are unchanged when a precision is multiplied by a positive constant. For unequal blocks $\Pi^{\mathcal P}=\diag(n_\alpha)^{1/2}\Phi^\top\Pi\Phi\diag(n_\alpha)^{1/2}$ is a congruence by a diagonal matrix that is not a multiple of the identity, which changes contrasts in general (\cref{chk:col-coarse-medium}). The orthonormal block precision $\Phi^\top\Pi\Phi$ can also leave the diffusive class: for a block of three agents on a path with unit couplings and a block of one agent coupled to an end of the path with coupling $\frac12$, at leak $0.01$, the first row sum of $\Phi^\top\Pi\Phi$ is $0.53/3-0.5/\sqrt3<0$.
\end{proof}
% src: H6-L1 proof; the conditional reading of (b) is derived here

\begin{proof}[Proof of \cref{prop:col-nesting}]
(a) Collective variables of $C_k$ are collective variables of $C$, so a variable that is not collective for $C$ is not collective for $C_k$. A member variable of $C_k$ in $\sS(C)$ reads such a variable and lies in $\sS(C_k)$; one in $\sA(C)$ drives such a variable and lies in $\sA(C_k)$ unless it reads outside $C_k$, in which case it lies in $\sS(C_k)$. For exactness, use \cref{thm:col-path-blanket}(b) with all other channels separate. A channel of a member of $C_1$ that superposes a message from a member of $C_2$, read only by members of $C$ and so internal to $C$, and a message from a non-member of $C$ reads two variables exterior to $C_1$, so it is no collider for $C_1$, but it is a collider for $C$. A channel of a member of $C_1$ that superposes messages from another member of $C_1$, internal to $C_1$, and from a member of $C_2$ is a collider for $C_1$, while for $C$ the channel and both sources are internal.

(b) Let $B(x,r)$ and $B(y,s)$ be closed $u$-balls with $r\le s$ and a common point $z$. For $w\in B(x,r)$, $u(w,y)\le\max\{u(w,x),u(x,z),u(z,y)\}\le s$, so $B(x,r)\subseteq B(y,s)$; concatenating chains shows that $u$ is an ultrametric. A set $C$ with $\mathcal U(C)>0$ is a ball: at any height $h\in[u_{\mathrm{in}}(C),\varpi_{\mathrm{out}}(C))$ its members are equivalent, and a chain from a member to a non-member contains a link across the boundary of weight at least $\varpi_{\mathrm{out}}(C)>h$. For several modalities, with $u_\vartheta=\max_cu_{(c)}/\vartheta_c$,
\[
\frac{u_{(c)}(i,k)}{\vartheta_c}\le\frac{\max\{u_{(c)}(i,j),u_{(c)}(j,k)\}}{\vartheta_c}\le\max\{u_\vartheta(i,j),u_\vartheta(j,k)\},
\]
and taking the maximum over $c$ shows that $u_\vartheta$ is an ultrametric, whose balls of radius one are the intersections $\bigcap_cB_{u_{(c)}}(i,\vartheta_c)$. For crossing, take three agents with $\varpi^{(1)}_{ij}$ and $\varpi^{(2)}_{jk}$ small and all other weights large: at a height between the small and the large weights, modality $1$ has the ultrametric class $\{i,j\}$ and modality $2$ the class $\{j,k\}$, and $\min_cu_{(c)}$ is small on $(i,j)$ and $(j,k)$ and large on $(i,k)$, violating the ultrametric inequality.

(c) The canonical partition is defined for channels of every modality, so a channel that reads or drives a non-member lies in the blanket whatever its modality. The modality-restricted partition places such a channel of another modality in the interior although it reads or drives an exterior variable, which violates $A_{\sI\sE}=0$ or $A_{\sE\sI}=0$. In the split partition the only arrows between interior and exterior are the within-member couplings named in the statement, which are direct interior-exterior entries of $A$, and \cref{thm:col-path-blanket}(a) applies.
\end{proof}
% src: L2-P11 proof, L5-P3 proof, L5-P4 proof, L2-P12 proof; the step from positive contrast to balls is written out here

\begin{proof}[Proof of \cref{thm:col-boundary}]
Part (a) combines \cref{thm:col-path-blanket}, \cref{prop:col-records-partition} and \cref{prop:col-two-hop}; a one-hop shell in which no shell agent reads both sides has a vanishing collider sum. Part (b) applies \cref{thm:col-thick-shell,thm:col-markov-bound} to the world-conditional law, which is Gaussian by \cref{prop:col-world-conditional}, with the constants of \cref{thm:col-markov-bound}(c) under time-invariant dynamics with local noise and row-sum contraction, and \cref{thm:col-diffusive-precision}(c) to the reversible and normal isotropic cases. Part (c) is \cref{prop:col-markov-vs-correlation}, and (e) is \cref{prop:col-selects-nothing} with \cref{def:col-collective}.

For (d), the bound by memory, the restoration by the filtered estimate and the reversible case are \cref{prop:col-memory}. That the instantaneous blanket can fail is shown by one agent whose sensory channel $s$ reads a world variable $e$ and whose internal state $i$ reads $s$: $X_{e,t+1}=\varphi X_{e,t}+\Xi^e_t$ with $0<\varphi<1$, $X_{s,t+1}=X_{e,t}+\Xi^s_t$ and $X_{i,t+1}=X_{s,t}+\Xi^i_t$, with unit noise variances, $\sI=\{i\}$, $\sB=\{s\}$, $\sE=\{e\}$. No variable reads both sides, so the path blanket holds by \cref{thm:col-path-blanket}. In the stationary law, with $\sigma^2=1/(1-\varphi^2)$, $\Cov(X_{e,t},X_{s,t})=\varphi\sigma^2$, $\Var(X_{s,t})=\sigma^2+1$, $\Cov(X_{i,t},X_{s,t})=\varphi\sigma^2$ and $\Cov(X_{i,t},X_{e,t})=\varphi^2\sigma^2$, since $X_{i,t}$ depends on the world through $X_{e,t-2}$. The partial covariance of $X_{i,t}$ and $X_{e,t}$ given $X_{s,t}$ is $\varphi^2\sigma^2-\varphi^2\sigma^4/(\sigma^2+1)=\varphi^2\sigma^2/(\sigma^2+1)>0$, so the instantaneous leak is positive; it is $0.039$ nats at $\varphi=0.9$ (\cref{chk:col-path}).
\end{proof}
% src: B1-T4 (assembled); L2-E1 for the example of (d), whose positivity is derived here in closed form
